\documentclass[11pt]{article}
\usepackage[a4paper,margin=22mm]{geometry}
\usepackage[T1]{fontenc}
\usepackage{amsmath,amssymb,amsthm,mathtools,bm}
\usepackage{graphicx,booktabs,enumitem,hyperref}
\usepackage[capitalize,noabbrev]{cleveref}
\newtheorem{theorem}{Theorem}
\newtheorem{lemma}{Lemma}
\newtheorem{proposition}{Proposition}
\newtheorem{corollary}{Corollary}
\theoremstyle{definition}
\newtheorem{definition}{Definition}
\newtheorem{assumption}{Assumption}
\theoremstyle{remark}
\newtheorem{remark}{Remark}
\crefname{assumption}{Assumption}{Assumptions}
\Crefname{assumption}{Assumptions}{Assumptions}
\newcommand{\E}{\mathbb E}
\newcommand{\R}{\mathbb R}
\newcommand{\one}{\mathbf 1}
\newcommand{\KL}{D_{\mathrm{KL}}}
\newcommand{\softmax}{\operatorname{softmax}}
\newcommand{\relu}{\operatorname{ReLU}}
\newcommand{\ip}[2]{\left\langle #1,#2\right\rangle}
\newcommand{\norm}[2][2]{\left\lVert #2\right\rVert_{#1}}
\newcommand{\define}{\triangleq}
\newcommand{\pd}[2]{\mathchoice{\frac{\partial #1}{\partial #2}}{\tfrac{\partial #1}{\partial #2}}{\partial #1/\partial #2}{\partial #1/\partial #2}}
\newcommand{\dpd}[2]{\frac{\partial #1}{\partial #2}}
\newcommand{\why}[1]{\text{\small(#1)}}
\newcommand{\ddt}[1]{\mathchoice{\frac{d#1}{dt}}{\tfrac{d#1}{dt}}{d#1/dt}{d#1/dt}}
% Shared mathematical notation. Change a symbol here, not in the section files.
% Auxiliary symbols are defined where used; there is no separate notation dictionary.
\newcommand{\Lout}{\mathcal L_{\mathrm{out}}}
\newcommand{\Lstep}{\mathcal L_{\mathrm{step}}}
\newcommand{\cred}{\Phi}
\newcommand{\outcredit}{\cred_{\mathrm{out}}}
\newcommand{\rulecredit}{\cred_{\mathrm{rule}}}
\newcommand{\thout}{\bm\theta^\star_{\mathrm{out}}}
\newcommand{\thstep}{\bm\theta^\star_{\mathrm{step}}}
\newcommand{\initscale}{\varepsilon}
\newcommand{\ballradius}{r_0}
\newcommand{\failprob}{\delta}
\newcommand{\stepsize}{\eta}
\newcommand{\mixrate}{\lambda}
\newcommand{\initball}{\mathcal B_{\initscale}}
\newcommand{\traintime}{\tau}
\newcommand{\basin}{\mathcal U}
\newcommand{\items}{\mathcal I}
\newcommand{\pimin}{\pi_{\min}}
\newcommand{\margin}{\gamma}
\newcommand{\execmargin}{\margin^\star}
\newcommand{\rulemargin}{\gamma_{\mathrm{rule}}}
\newcommand{\embconst}{c_{\mathrm{emb}}}
\newcommand{\entryidx}{j_0}
\newcommand{\Ktokens}{K_{\mathrm{tokens}}}
\newcommand{\attconst}{C_{\mathrm{att}}}
\newcommand{\Jz}{\mathchoice{\frac{\partial\bm z_{\bm\theta}}{\partial\bm\theta}}{\tfrac{\partial\bm z_{\bm\theta}}{\partial\bm\theta}}{\partial\bm z_{\bm\theta}/\partial\bm\theta}{\partial\bm z_{\bm\theta}/\partial\bm\theta}}
\newcommand{\term}{\mathsf t}
\newcommand{\proofstep}[1]{\par\smallskip\noindent\textbf{#1}\enspace}
\newsavebox{\fitbox}
\newcommand{\fitwidth}[1]{\sbox{\fitbox}{#1}\ifdim\wd\fitbox>\linewidth\resizebox{\linewidth}{!}{\usebox{\fitbox}}\else\usebox{\fitbox}\fi}
\title{Counting Transformers: Separate Direct-Proof Draft}
\author{Working mathematical draft}
\date{10 October 2026}
\setlength{\emergencystretch}{2em}
\begin{document}
\maketitle
\noindent This document is separate from the manuscript. It preserves its four headline theorem conclusions. The outcome loss floor is proved by conditioning on the letters used by each network term, without an interaction decomposition. A short conditional-projection construction retains the degreewise credit results. The network expansion and the binary joint-flow proof remain necessary; the latter is reproduced in full with a corrected auxiliary numerical constant. No general-$K$ joint-convergence theorem is claimed.

\paragraph{Reading order.} The new argument is in Sections~\ref{sec:conditioning} and~\ref{sec:direct-outcome}; Section~\ref{sec:retain-credit} retains the mechanistic result. The architecture estimates and joint-flow proof are included so that no proof lemma needs to be imported from the manuscript. Generic constants can enlarge from one line to the next and depend only on the fixed architecture depth, number of heads, $K$, $a_0$ and $c_{\mathrm{emb}}$.

\section{Setup and preserved statements}\label{sec:setting}
\begin{definition}[Counting modulo $K$]
\label{def:task}
A prompt has a start count $s_0$, uniform on $\mathbb Z_K$, and letters $w_1,\ldots,w_n$ drawn i.i.d.\ from an alphabet
in which a fixed query letter $w^\star$ has probability $p\in(0,1)$. With $b_t=\mathbf 1[w_t=w^\star]$,
\begin{equation}
 s_t=s_{t-1}+b_t\pmod K,\qquad t=1,\ldots,n,
 \label{eq:tracking}
\end{equation}
and the answer is $y=s_n$.
\end{definition}

\begin{definition}[Contraction coefficient]
\label{def:lambda}
For counting modulo $K$,
\begin{equation}
 \mixrate=\max_{1\le j<K}\bigl|1-p+p\,e^{2\pi\mathrm ij/K}\bigr|=\bigl|1-p+p\,e^{2\pi\mathrm i/K}\bigr|.
 \label{eq:lambda}
\end{equation}
\end{definition}


\begin{definition}[Outcome and step losses]
\label{def:losses}
Let $\bm z_t(\bm\theta)\in\R^K$ be the state logits at the position that predicts $s_t$, and $\bm z_{\bm\theta}$
those at the answer position. The \emph{outcome loss} is $\Lout(\bm\theta)=\E[-\log\softmax(\bm z_{\bm\theta})_y]$, and $\bm y=\bm e_y$ is the one-hot answer.
Write $\ell_t(\bm\theta)=\E[-\log\softmax(\bm z_t(\bm\theta))_{s_t}]$. The first state prediction is a boundary case: the previous
count $s_0$ sits in the prompt rather than at the predicting position. The \emph{step loss} is therefore the teacher-forced
transition loss
\[
 \Lstep(\bm\theta)=\frac{1}{n-1}\sum_{t=2}^n\ell_t(\bm\theta),
\]
and $\Lstep^{\mathrm{all}}(\bm\theta)=\tfrac{1}{n}\sum_{t=1}^n\ell_t(\bm\theta)$ also includes $t=1$.
\end{definition}

\begin{definition}[Degree]
\label{def:degree}
Fix the start count. The orthogonal conditional-projection component (constructed below) $f_S$ of a square-integrable function $f$ of the letters has
\emph{degree} $|S|$, and $(f)_r=\sum_{|S|=r}f_S$. Vector- and matrix-valued functions are decomposed entrywise.
\end{definition}
For centred match bits, a sum has degree $1$ and a product of two distinct bits has degree $2$.
Degree counts input letters, not powers of parameters. Write $A_r:=\E[\|(\bm y-\bm u)_r\|^2]$ and $J_r(\bm\theta):=\E[\|(\Jz)_r\|_F^2]$ as in \cref{eq:matching}; the
\emph{interaction credit} is the credit at degrees $r\ge2$, and $\outcredit=\cred_{\bm y-\bm u}$. With the state-only predictor
$\bm q_s$ and the rule signal $\bm g_t$ of \cref{sec:setting}, $\rulecredit=\E_t\E[\langle\bm z_t,\bm g_t\rangle]$ with $t$ uniform on $\{2,\ldots,n\}$.

\paragraph{Architecture and initialization.}

\begin{assumption}[Architecture]
\label{ass:arch}
A causal Transformer with two blocks of softmax heads and ReLU MLPs, learned token and position
embeddings, a learned relative position bias $\beta_{i-j}$ added to every attention score and initialized at zero, no
LayerNorm, any widths (residual $d_{\mathrm{res}}$, head $d_{\mathrm h}$, MLP $w$) and a linear readout $\mathbf V$;
every logit vector is $\bm z=\mathbf V\bm h$ with $\bm h$ independent of $\mathbf V$. A head may have a fixed, input-independent causal submask;
ordinary attention uses the full prefix.
\end{assumption}

\begin{assumption}[Initialization]
\label{ass:init}
\emph{Scale.} $\bm\theta_0=\initscale\tilde{\bm\theta}$ with $\initscale\le\initscale_0:=c\min(1,\frac{1}{N_b})$, where every weight matrix and embedding block of $\tilde{\bm\theta}$ has operator norm
at most $a_0$ and the value-token embeddings have smallest singular value at least $\embconst>0$.

\emph{Biases.} Every MLP bias satisfies $|\tilde b_m|\ge1$, and each MLP bias vector has norm at most $N_b$.

\emph{Neighbourhood.} We work on
$\initball=\{\bm\theta:\|\bm\theta-\bm\theta_0\|\le\ballradius\initscale,\ \beta=0\}$ with $\ballradius=\min(\frac{1}{8},\frac{\embconst}{4})$, where $\beta$ is the
relative position bias.
\end{assumption}

\begin{assumption}[Task]
\label{ass:task}
The task is counting modulo $K\ge2$ (\cref{def:task}), so $\mixrate<1$.
\end{assumption}

\paragraph{Constants and widths.} On $\initball$, no ReLU unit changes sign (\cref{app:network}). Constants $c,C>0$ depend only on $a_0$, $\embconst$, $K$
and the numbers of blocks and heads, never on the widths, $n$, $\initscale$ or the time. Fan-in-scaled initialization
satisfies \cref{ass:init} w.h.p.\ at every width, with $a_0$ bounded in terms of the largest block aspect ratio,
$\embconst\ge\tfrac{1}{2}$ and $N_b\le2\sqrt w$, provided the MLP biases are drawn with $|\tilde b_m|\in[1,2]$, which default
initializations (zero or small biases) do not do; at the executor width the aspect ratios are bounded, so the theorems
hold there with $n$-independent constants for $\initscale\le c\min(\mixrate,n^{-1/2})$ (the initialization discussion above).

\paragraph{Parameters and training.}

\begin{assumption}[Training]
\label{ass:gd}
Full-batch gradient flow $\tfrac{d\bm\theta}{dt}=-\mathbf P\nabla_{\bm\theta}\mathcal L(\bm\theta)$ on $\Lout$ or $\Lstep$, where $\mathbf P$ is the orthogonal projection onto
the trainable parameter blocks, with the relative
position bias held at its zero initialization. \Cref{thm:outcome-end} drops this restriction and lets the bias train (\cref{lem:bias-uniform}).
\end{assumption}

\label{rem:norms}Parameters carry the Euclidean norm of the vectorized parameters (Frobenius on
each block). In \cref{thm:main,thm:separation} the relative bias is fixed at zero, which is what $\beta=0$ in $\initball$ and
\cref{ass:gd} encode, and $\bm\theta$ denotes the parameters other than $\beta$: the Jacobian $\Jz$, the gradients $\nabla_{\bm\theta}\mathcal L$ and
$\nabla_{\bm\theta}\cred$, and the degree spectra $J_r$ are taken with respect to these parameters only. The step-credit bounds use uniform
zero-score attention through \cref{prop:network}(iv) and do not cover a trained bias. \Cref{thm:outcome-end} allows the bias to be
trained: its proof uses only the logit bounds, which hold for every $\beta$ (\cref{lem:bias-uniform}). The runs of the manuscript routing experiments also train it; the gate head of \cref{thm:joint} has no relative bias. The in-regime runs of the manuscript experiments have no relative bias.



\paragraph{Conventions.} (A1)--(A4) denote \cref{ass:arch,ass:init,ass:task,ass:gd}; (A3$'$) means a uniform start count. A query is fixed throughout each conditional argument, with query-match probability $p$. The target $\bm y$ denotes the one-hot answer. All local estimates also hold on the doubled ball $\initball'$ of relative radius $2r_0$. Since $2r_0\le\min(1/4,c_{\mathrm{emb}}/2)$, biases there have magnitude at least $\varepsilon/2$, and the smallest singular value of the count-embedding block is at least $\varepsilon c_{\mathrm{emb}}/2$. We use $0<\lambda<1$ whenever inverse powers of $\lambda$ occur. At $\lambda=0$, the hypotheses involving $\varepsilon\le c\lambda$ admit no positive initialization scale.

\begin{definition}[Credit]\label{def:credit}
For a signal $g$, define $\cred_g(\theta)=\mathbb E\langle z_\theta,g\rangle$. Its credit is $\|\nabla_\theta\cred_g\|$. Interaction credit means the norm of the sum of matching-degree contributions at degrees at least two; the proof below also bounds every single such degree separately. The relative bias is excluded from the credit Jacobian and fixed at zero in Theorems 1 and 2.
\end{definition}
\begin{theorem}[Outcome interaction credit]
\label{thm:main}
Under \cref{ass:arch,ass:init,ass:task}, with the start count fixed, for all $n\ge2$, $\initscale^2\le c\mixrate$ and
$\bm\theta\in\initball$, the outcome credit at degrees $r\ge2$ is at most $C\initscale^6\mixrate^{n-3}$.
\end{theorem}

\begin{theorem}[Step credit]
\label{thm:separation}
Under \cref{ass:arch,ass:init,ass:task} with a uniform start, for all $n\ge2$ and $\bm\theta\in\initball$:
\begin{enumerate}[label=(\roman*),leftmargin=*,itemsep=0pt,topsep=2pt]
\item \emph{(prior-fitting gradient)} if $\initscale\le c\mixrate$, then $\|\nabla_{\bm\theta}\Lstep(\bm\theta)\|\ge c\,\initscale\mixrate$ while
$\|\nabla_{\bm\theta}\Lout(\bm\theta)\|\le C\initscale(\initscale^2+\mixrate^n)$;
\item \emph{(rule credit)} if $\initscale\le c$ and attention uses the full causal prefix, then $\|\nabla_{\bm\theta}\rulecredit(\bm\theta)\|\le \frac{C\initscale^6}{n}$; if
moreover $p<\tfrac{1}{2}$, the degree-$6$ polynomial $P_n$ is not identically zero, with an explicit count$\times$letter monomial of
coefficient of order $\frac{1}{n\sqrt{d_{\mathrm h}}}$
(\cref{lem:leading}).
\end{enumerate}
\end{theorem}

\begin{theorem}[Exponential-time outcome plateau near small initialization]
\label{thm:outcome-end}
Under \cref{ass:arch,ass:init,ass:task} with a zero initial readout, let $P(t)$ be, at each time, the orthogonal projection onto a set
of trainable parameter blocks, piecewise constant in $t$; the blocks may include the relative bias. Consider gradient flow
$\tfrac{d\bm\theta}{dt}=-\mathbf P(t)\nabla_{\bm\theta}\Lout(\bm\theta)$; full joint flow ($\mathbf P\equiv\mathbf I$) and training only the blocks of \cref{thm:joint} are special cases.
For all $n\ge2$, $\initscale\le c\mixrate$ and $0\le t\le c\mixrate^{-n}$, with constants that do not depend on the schedule:
\begin{enumerate}[label=(\roman*),leftmargin=*,itemsep=0pt,topsep=2pt]
\item \emph{(locality)} $\|\bm\theta(t)-\bm\theta_0\|\le\ballradius\initscale$, the relative bias included in the norm;
\item \emph{(loss)} $\Lout(\bm\theta(t))\ge\log K-C\initscale^2\mixrate^n$;
\item \emph{(correct-answer probability)} $\E\bigl[\softmax(\bm z_{\bm\theta(t)})_y\bigr]\le\frac{1}{K}+C\initscale\mixrate^{n/2}$.
\end{enumerate}
\end{theorem}

\begin{theorem}[Joint gradient flow learns the transition]
\label{thm:joint}
Count modulo~$2$ with $n\ge2$ letters and a query of probability $p\in(0,1)$, $p\neq\tfrac{1}{2}$ (so $\mixrate=|1-2p|>0$), and let
$h^\star=|\log\tfrac{1-p}{p}|$ and $\mathrm{gap}\in(0,1]$. There is $\initscale^\star>0$, explicit in $\mathrm{gap}$, $p$ and $d_{\mathrm h}$ (\cref{app:J-constants}), such that for
$0<\initscale\le\initscale^\star$ and every initialization with $|D|\ge\mathrm{gap}\,\initscale^2$, where $D=q(0)^2-k(0)^2$, joint gradient flow on $\Lstep$ gives every prediction positive margin at
every time $t\ge T$, and $\Lstep\to0$, where, with an absolute constant $C$,
\[
 T\le C\Bigl(\frac{\sqrt{d_{\mathrm h}}}{p(1-p)\,h^\star\mathrm{gap}\,\initscale^4}+\frac{d_{\mathrm h}}{\min(p,1-p)^2\,\initscale^{12}}\Bigr).
\]
$T$ depends on $n$ only through $d_{\mathrm h}$ (any $d_{\mathrm h}\ge1$ is allowed). If $q(0)k(0)(1-2p)<0$, the condition on $D$ and the
first term are not needed. If $\frac{|q(0)|}{\initscale}$ and $\frac{|k(0)|}{\initscale}$ are independent and uniform on $[\tfrac{1}{2},2]$, $\Pr(|D|<\mathrm{gap}\,\initscale^2)\le\tfrac{4}{3}\mathrm{gap}$.
\end{theorem}

\section{Direct contraction by conditioning}\label{sec:conditioning}
Let $G=\bm y-\bm u$, $\bm u=K^{-1}\bm1$, $\alpha=e^{2\pi\mathrm i/K}$ and $\rho_j=1-p+p\alpha^j$. Fix $s_0$ and the query. For $S\subseteq[n]$, let $\mathcal F_S=\sigma(w_t:t\in S)$.

\begin{lemma}[Conditional target contraction]\label{lem:conditional-target}
For every $S\subseteq[n]$,
\[
 \|\mathbb E[G\mid\mathcal F_S]\|^2
 =\frac1K\sum_{j=1}^{K-1}|\rho_j|^{2(n-|S|)}
 \le (1-K^{-1})\lambda^{2(n-|S|)}.
\]
For every square-integrable $\mathcal F_S$-measurable vector $H$, and every matrix-valued such $H$ with compatible dimensions,
\[
 |\mathbb E\langle H,G\rangle|\le\lambda^{n-|S|}\|H\|_{L^2},\qquad
 \|\mathbb E[H^\top G]\|\le\lambda^{n-|S|}\|H\|_{L^2(F)}.
\]
\end{lemma}
\begin{proof}
The nonconstant Fourier coefficient of the count distribution at frequency $j$ is the conditional expectation of $\alpha^{js_n}$. Independence gives
\[
 \mathbb E[\alpha^{js_n}\mid\mathcal F_S]
 =\alpha^{js_0}\rho_j^{n-|S|}\prod_{t\in S}\alpha^{jb_t}.
\]
Every factor in the product has modulus one. Parseval for the $K$-point discrete Fourier transform, with its $1/K$ normalization on probability vectors, gives the asserted identity and bound. Since $H$ is measurable on $\mathcal F_S$, the tower property gives $\mathbb E\langle H,G\rangle=\mathbb E\langle H,\mathbb E[G\mid\mathcal F_S]\rangle$. Cauchy--Schwarz proves the scalar bound. Testing the matrix expression against an arbitrary unit vector proves the vector bound. The factor $\sqrt{1-1/K}\le1$ was dropped.
\end{proof}

\paragraph{Which objects require conditioning.} The argument applies to a network term supported on $S$, even when that term has nonzero conditional means. It therefore permits summing the ordinary attention-expansion terms directly. No decomposition by orthogonal interaction degree is used in the outcome proof below.

\section{Direct local credit and the outcome plateau}\label{sec:direct-outcome}
For fixed parameters, introduce an auxiliary variable $u$: multiply every content attention score, in both blocks, by $u$, leaving any positional relative bias unchanged. The second-block residuals and hence its content scores depend on $u$. This is the composed network, rather than an expansion with frozen second-block scores. Define
\begin{equation}
 z_\theta=\sum_{k=0}^N z^{[k]}+R_N,\qquad
 z^{[k]}=\frac1{k!}\frac{d^kz(u)}{du^k}\bigg|_{u=0},\qquad
 R_N=\int_0^1\frac{(1-u)^N}{N!}\frac{d^{N+1}z(u)}{du^{N+1}}\,du.
 \label{eq:score-split}
\end{equation}
For credit Jacobians let $J^{[k]}=\partial_\theta z^{[k]}$ and $J_R=\partial_\theta R_N$. The auxiliary Taylor order is not the parameter-polynomial degree.

\begin{lemma}[Support and absolute coefficient bounds]\label{lem:direct-support}
On the doubled initialization ball, with relative bias fixed at zero, write every positive-order coefficient as a finite sum of terms $z^{[k]}=\sum_\nu H_{k,\nu}$, each depending on at most $2k+1$ letters. There are constants $C,B$ independent of $n$, widths and $k$ such that for $k\ge1$,
\[
 \sum_\nu\sup\|H_{k,\nu}\|\le C\varepsilon^3(B\varepsilon^4)^k,\qquad
 \sum_\nu\sup\|\partial_\theta H_{k,\nu}\|_F\le C\varepsilon^2(B\varepsilon^4)^k.
\]
The remainder satisfies $\sup\|R_N\|\le C\varepsilon^3(B\varepsilon^4)^{N+1}$ and $\sup\|J_R\|_F\le C\varepsilon^2(B\varepsilon^4)^{N+1}$. At the answer position the zero-order coefficient is a fixed-input term of size $C\varepsilon^2$ and tangent size $C\varepsilon$, plus a sum of one-letter terms with absolute sizes at most $C\varepsilon^3$ and tangent sizes at most $C\varepsilon^2$.

The logit assertions hold uniformly for arbitrary relative bias $\beta$, provided the other parameters stay in the doubled ball.
\end{lemma}
\begin{proof}
The complete product, softmax, two-block and marked-site calculations are given in \cref{prop:network}. Its absolute coefficient bounds give the first assertion before any cancellation; its input-count calculation gives support at most $2k+1$. For a map into $\mathbb R^K$, its parameter Jacobian has Frobenius norm at most $\sqrt K$ times its operator norm. Thus the directional tangent bounds there give the stated Frobenius bounds, with a fixed $\sqrt K$ absorbed in $C$. The integral in \eqref{eq:score-split} has normalized weight exactly matching the next normalized derivative, proving the remainder estimate.

At zero scores the attention weights are input independent. The direct answer-position embedding and every MLP offset depend on no random letters (conditional on the start and query). All transported letter terms contain the readout, a value matrix and a letter embedding. Their sizes are $O(\varepsilon^3)$, their tangents $O(\varepsilon^2)$, and their coefficients sum absolutely to a constant because all attention rows are probability vectors. This proves the zero-order assertion.

For arbitrary $\beta$, replace the zero-score row $m_i^{-1}\bm1$ by $\operatorname{softmax}(\beta_{i-\cdot})$. Every such row has nonnegative entries summing to one. The derivative identity and normalized softmax derivative estimate in \eqref{eq:softmax-avg} hold at every real score vector. At $u=0$ the coefficients are weighted averages with those fixed probabilities; the absolute sum estimate is unchanged. Biases add no input sites and do not differentiate with respect to $u$, so the support induction is unchanged as well. The residual bounds imply $|(W_{\rm in}h)_m|\le C\varepsilon^2<|b_m|$, independently of $\beta$, keeping the activation pattern fixed. Only the marked-site $1/m_i$ estimate in \cref{prop:network}(iv) requires $\beta=0$; it is not used for the logit assertions here.
\end{proof}

\begin{lemma}[Direct credit summation]\label{lem:score-sum}
Let $0<\lambda<1$, $N=\lceil(n-1)/2\rceil$ and $B\varepsilon^4\le\lambda^2/2$. Under the local assumptions,
\begin{align*}
 |\outcredit(\theta)|&\le C(\varepsilon^2\lambda^n+\varepsilon^3\lambda^{n-1}+\varepsilon^7\lambda^{n-3}),\\
 \|\nabla_\theta\outcredit(\theta)\|&\le C(\varepsilon\lambda^n+\varepsilon^2\lambda^{n-1}+\varepsilon^6\lambda^{n-3}).
\end{align*}
The first bound is uniform in arbitrary relative bias. In the second bound the bias is fixed at zero and excluded from the parameter Jacobian. If also $\varepsilon\le\lambda$, these bounds reduce to
\begin{equation}
 |\outcredit(\theta)|\le C\varepsilon^2\lambda^n,\qquad
 \|\nabla_\theta\outcredit(\theta)\|\le C\varepsilon\lambda^n.
 \label{eq:thm1-full}
\end{equation}
They remain valid after averaging over the start count.
\end{lemma}
\begin{proof}
Apply \cref{lem:conditional-target} directly to each coefficient term in \cref{lem:direct-support}. The zero-order bounds give $C(\varepsilon^2\lambda^n+\varepsilon^3\lambda^{n-1})$ for the logit credit and $C(\varepsilon\lambda^n+\varepsilon^2\lambda^{n-1})$ for its gradient. For $k\ge1$, every support size is at most $2k+1$, so
\[
 |\mathbb E\langle z^{[k]},G\rangle|
 \le C\varepsilon^3(B\varepsilon^4)^k\lambda^{n-2k-1}
 = CB\varepsilon^7\lambda^{n-3}
 \left(\frac{B\varepsilon^4}{\lambda^2}\right)^{k-1}.
\]
If $2k+1>n$, the displayed power of $\lambda$ is larger than one and still gives a valid upper bound. For the remainder use $\|G\|\le1$, hence
\[
 |\mathbb E\langle R_N,G\rangle|
 \le C\varepsilon^3(B\varepsilon^4)^{N+1}
 =CB\varepsilon^7\lambda^{2N}
 \left(\frac{B\varepsilon^4}{\lambda^2}\right)^N
 \le CB\varepsilon^7\lambda^{n-3}\,2^{-N},
\]
since $2N\ge n-1$. Summing a geometric series proves the first bound. Apply the same conditional estimate to the Jacobian terms, with one fewer power of $\varepsilon$, for the second bound. If $\varepsilon\le\lambda$ and $\varepsilon\le1$, then $\varepsilon/\lambda\le1$ and $\varepsilon^5/\lambda^3\le\varepsilon^2\le1$, which proves \eqref{eq:thm1-full}. All bounds were uniform in the fixed start, so averaging preserves them.
\end{proof}

\begin{lemma}[Uniform loss floor and smoothness]\label{lem:bias-uniform}
On the set $\mathcal B^{\rm full}_\varepsilon=\{(\vartheta,\beta):\|\vartheta-\vartheta_0\|\le2r_0\varepsilon\}$, with $\beta$ unrestricted, if $\varepsilon\le c\lambda$ then
\[
 |\outcredit|\le C_2\varepsilon^2\lambda^n.
\]
The loss has gradient norm at most $C_1\varepsilon$ and directional second derivatives at most $C_1$, for the full parameter norm including $\beta$.
\end{lemma}
\begin{proof}
Choose $c$ small enough that $\varepsilon\le\lambda$ and $B\varepsilon^4\le\lambda^2/2$. The credit estimate follows from \cref{lem:score-sum}. For smoothness, a unit perturbation of the relative bias changes a row of scores by a vector of supremum norm at most one. By \eqref{eq:softmax-avg}, its attention row has bounded first and second derivatives in $\ell_1$, uniformly in its number of visible sites and in $\beta$. All parameter matrices and input embeddings have directional profiles $(C\varepsilon,C,0)$. A bias vector has norm at most $C$ and profile $(C,1,0)$; it enters through $o=W_{\rm out}D b$, whose profile is $(C\varepsilon,C,C)$. Products have profiles
\[
 (A,A',A'')(B,B',B'')=(AB,A'B+AB',A''B+2A'B'+AB'').
\]
Induct over the two blocks. Residuals have profile $(C\varepsilon,C,C)$; projected queries, keys and values have profile $(C\varepsilon^2,C\varepsilon,C)$; content scores have profile $(C\varepsilon^4,C\varepsilon^3,C\varepsilon^2)$. Including the relative bias leaves attention rows with profile $(1,C,C)$ in $\ell_1$. A head output has profile $(C\varepsilon^2,C\varepsilon,C)$ and the final logits have profile $(C\varepsilon^2,C\varepsilon,C)$. The fixed activation pattern makes every intervening MLP affine. For a single-example cross-entropy, $p=\operatorname{softmax}(z)$,
\[
 \ell'=\langle p-y,z'\rangle,\qquad
 \ell''=z'^\top(\operatorname{diag}p-pp^\top)z'+\langle p-y,z''\rangle.
\]
Since $\|p-y\|\le\sqrt2$ and $\|\operatorname{diag}p-pp^\top\|\le1/2$, these are bounded by $C\varepsilon$ and $C$. Taking expectations and testing arbitrary unit directions proves the assertions. Every estimate used operator norms or probability-row $\ell_1$ norms, so the constants have the claimed independence.
\end{proof}
\label{lem:end-smooth}

\paragraph{Loss identity.}\label{app:end}
For $p_\theta=\operatorname{softmax}(z_\theta)$, direct algebra gives
\begin{equation}
 L_{\rm out}(\theta)=\log K+\mathbb E D_{\rm KL}(u\|p_\theta)-\outcredit(\theta).
 \label{eq:lo-identity}
\end{equation}
Indeed $D_{\rm KL}(u\|p)=\log\sum_a e^{z_a}-\log K-\langle u,z\rangle$ and the cross-entropy is $\log\sum_a e^{z_a}-\langle y,z\rangle$.

\begin{proof}[Complete proof of \cref{thm:outcome-end}]
Set $B_n=C_2\varepsilon^2\lambda^n$. On the doubled local set, \eqref{eq:lo-identity} and nonnegativity of KL give
\begin{equation}L_{\rm out}(\theta)\ge\log K-B_n.\label{eq:lo-floor}\end{equation}
At zero initial readout, $L_{\rm out}(\theta_0)=\log K$. Let $\tau$ be the first exit time from the closed full parameter ball of radius $r_0\varepsilon$, including the relative bias. For piecewise constant block projections (with finitely many switches on each finite time interval), local smoothness gives a unique trajectory between switches and a continuous trajectory across them. Almost everywhere before exit,
\[
 \frac{dL_{\rm out}}{dt}=-\|P(t)\nabla L_{\rm out}\|^2=-\|\dot\theta\|^2.
\]
Consequently, for $t\le\tau$,
\[
 \int_0^t\|\dot\theta(s)\|^2ds\le B_n,\qquad
 \|\theta(t)-\theta_0\|^2\le t\int_0^t\|\dot\theta(s)\|^2ds\le tB_n.
\]
If exit occurs, continuity gives $r_0^2\varepsilon^2\le\tau B_n$, hence $\tau\ge r_0^2 C_2^{-1}\lambda^{-n}$. If it never occurs, this lower bound is automatic. Taking $c\le r_0^2/(2C_2)$ proves locality on the claimed interval and proves the loss floor there. The argument also holds for measurable block projections using an absolutely continuous flow, since the same energy identity holds almost everywhere; thus permitting accumulation of switching times does not alter the bound.

Finally descent and the identity give $\mathbb E D_{\rm KL}(u\|p_t)\le\outcredit(\theta(t))\le B_n$. For every example, $p_{t,y}-1/K\le\operatorname{TV}(p_t,u)$. Pinsker's inequality and Jensen therefore give
\[
 \mathbb E[p_{t,y}]\le K^{-1}+\sqrt{B_n/2}
 \le K^{-1}+C\varepsilon\lambda^{n/2}.
\]
For completeness, Pinsker's inequality here follows by grouping the coordinates where $u_a\ge p_a$, applying the log-sum inequality to the two groups, and using $D_{\rm KL}((x,1-x)\|(q,1-q))\ge2(x-q)^2$. The last inequality follows by differentiating twice in $x$: its second derivative is $1/[x(1-x)]\ge4$ and it and its first derivative vanish at $x=q$. All time constants are independent of the block schedule.
\end{proof}

\section{Architecture estimates with complete derivation}
\label{app:network}

\begin{proposition}[Bounds for the score expansion]
\label{prop:network}
Assume (A1)--(A2). Uniformly on $\initball$, at every position and input, with $\bm x_i$ the token-plus-position vector,
$\partial_vF=\frac{d}{ds}F(\bm\theta+sv)|_{s=0}$, and $\|F\|_{\mathrm{abs}}$ the sum of the monomial norm bounds before cancellation:
\begin{enumerate}[label=(\roman*),leftmargin=*,itemsep=1pt,topsep=2pt]
\item with input-independent offsets $\bm O^1=\bm o_1$, $\bm O^2=\bm o_1+\bm o_2$: $\bm h_i^\ell=\bm x_i+\bm O^\ell+O(\initscale^2)$, $|\psi_{il}|\le C\initscale^4$,
$\|\bm z_i\|\le C\initscale^2$, $\|\partial_v\bm z_i\|\le C\initscale\|v\|$;
\item for $k\ge1$, $u\in[0,1]$ and $\|v\|=1$,
\[
 \frac{1}{k!}\Bigl\|\frac{d^k\bm z_i}{du^k}\Bigr\|\le C\initscale^3(C_0\initscale^4)^k,\qquad
 \frac{1}{k!}\Bigl\|\partial_v\frac{d^k\bm z_i}{du^k}\Bigr\|\le C\initscale^2(C_0\initscale^4)^k;
\]
the same holds for the absolute sums at $u=0$, for $\bm h^2_i$ with one power of $\initscale$ fewer, and for the Taylor remainder with $k=k_0+1$;
\item the zero-score network is affine in the inputs; its transported logit terms have tangent $O(\initscale^2)$, and direct-input and
input-free terms have tangent $O(\initscale)$; an order-$k$ coefficient contains at most $2k+1$ inputs; a one-score logit term has size
$O(\initscale^7)$ and tangent $O(\initscale^6)$;
\item with full-prefix attention, the terms of the order-$k$ coefficients at $u=0$ that contain an earlier input $\bm x_l$, $l<i$, obey the
bounds of (ii) multiplied by
\begin{equation}
 \kappa_{l,i}=\frac{1}{m_i}\Bigl(1+\sum_{l<j\le i}\frac{1}{m_j}\Bigr),\qquad m_i=i+1,
 \label{eq:dilution}
\end{equation}
after enlarging $C_0$.
\end{enumerate}
The constants do not depend on $n$, the widths or $k$.
\end{proposition}

\begin{proof}
For a network quantity $F$ let $F^{[k]}(u)=\frac{1}{k!}\frac{d^kF}{du^k}(u)$, so that $\bm z^{[k]}=\bm z^{[k]}(0)$, and write $\|F^{[k]}\|$ for its supremum over
inputs, positions and $u\in[0,1]$ (the $\ell_1$ norm for attention rows).

\emph{(i).} $|b_m|\ge\frac{\initscale}{2}>C\initscale^2\ge|(\mathbf W_{\mathrm{in}}\bm h)_m|$, so one activation matrix $\mathbf D_{\mathrm{act}}$ serves every input and every $u$, and
\[
 \bm h_i^\ell=\mathbf M_\ell\bigl(\bm h_i^{\ell-1}+\bm a_i^\ell\bigr)+\bm o_\ell,\qquad \mathbf M_\ell=\mathbf I+\mathbf W^\ell_{\mathrm{out}}\mathbf D_{\mathrm{act}}\mathbf W^\ell_{\mathrm{in}},\qquad \bm o_\ell=\mathbf W^\ell_{\mathrm{out}}\mathbf D_{\mathrm{act}}\bm b^\ell,
\]
with $\bm a_i^\ell=\sum_{\mathrm{heads}}\sum_{j\le i}\omega_{ij}\mathbf W_V\bm h^{\ell-1}_j$ and $\omega_{i\cdot}=\softmax(u\psi_{i\cdot})$. Then
\[
 \|\mathbf M_\ell-\mathbf I\|\le C\initscale^2,\qquad \|\bm o_\ell\|\le2a_0N_b\initscale^2\le a_0\initscale,\qquad \|\partial_v\bm o_\ell\|\le(a_0+2N_b)\initscale\le C,
\]
and $\|\omega_{i\cdot}\|_1=1$. Induction over the two blocks gives (i).

\emph{Softmax calculus.} For $\omega=\softmax(s)$, $\E_\omega[f]=\sum_j\omega_jf_j$ and a direction $\bm\nu$,
\[
 \tfrac{d}{dt}\omega_j(s+t\bm\nu)\big|_0=\omega_j(\nu_j-\E_\omega[\bm\nu]),\qquad
 \tfrac{d}{dt}\E_{\omega(s+t\bm\nu)}[f]\big|_0=\E_\omega[\bm\nu f]-\E_\omega[\bm\nu]\E_\omega[f].
\]
After $k$ derivatives each term has at most $k+1$ averages and each direction once; the absolute coefficient sums obey $a_{k+1}\le2(k+1)a_k$,
$a_0=1$. By duality,
\begin{equation}
 \bigl\|D^k\omega(s)[\bm\nu_1,\ldots,\bm\nu_k]\bigr\|_1\le2^kk!\prod_a\|\bm\nu_a\|_\infty .
 \label{eq:softmax-avg}
\end{equation}
Normalized derivatives of products are convolutions; those of compositions sum over the $\binom{k-1}{j-1}$ compositions of $k$ into $j$
positive parts. Hence
\[
 \tfrac{1}{j!}\|F^{(j)}\|\le Ac^j,\ \ \|h^{[k]}\|\le ab^k\quad\implies\quad\|(F\circ h)^{[k]}\|\le Aca\,b^k(1+ca)^{k-1}.
\]

\emph{(ii), values.} With $\|\bm\omega^{\ell,[0]}\|_1=1$, $\|\bm h^{0,[k]}\|=0$ for $k\ge1$, and $a_r\ge1$ in the compositions,
\begin{alignat*}{2}
 \|\psi^{\ell,[k]}\|&\le C\initscale^2\textstyle\sum_{a+b=k}\|\bm h^{\ell-1,[a]}\|\,\|\bm h^{\ell-1,[b]}\| &\qquad& \why{score is a product}\\
 \|(u\psi^\ell)^{[k]}\|&\le\|\psi^{\ell,[k]}\|+\|\psi^{\ell,[k-1]}\| && \why{$(u\psi)^{[k]}=u\psi^{[k]}+\psi^{[k-1]}$}\\
 \|\bm\omega^{\ell,[k]}\|_1&\le\textstyle\sum_{j=1}^k2^j\sum_{a_1+\cdots+a_j=k}\prod_r\|(u\psi^\ell)^{[a_r]}\| && \why{composition, \cref{eq:softmax-avg}}\\
 \|\bm a^{\ell,[k]}\|&\le C\initscale\textstyle\sum_{a+c=k}\|\bm\omega^{\ell,[a]}\|_1\,\|\bm h^{\ell-1,[c]}\| && \why{head output is a product}\\
 \|\bm h^{\ell,[k]}\|&\le C\bigl(\|\bm h^{\ell-1,[k]}\|+\|\bm a^{\ell,[k]}\|\bigr) && \why{affine block}\\
 \|\bm z^{[k]}\|&\le C\initscale\,\|\bm h^{2,[k]}\| && \why{readout}
\end{alignat*}
With $b=C\initscale^4$, block $1$ gives $\|\bm\omega^{1,[k]}\|_1\le b^k$ and $\|\bm a^{1,[k]}\|+\|\bm h^{1,[k]}\|\le C\initscale^2b^k$. Then $\|\psi^{2,[0]}\|\le C\initscale^4$,
$\|\psi^{2,[k]}\|\le C\initscale^5(k+1)b^k$ for $k\ge1$, and $\|(u\psi^2)^{[k]}\|\le C(2b)^k$, so
\begin{gather*}
 \|\bm\omega^{2,[k]}\|_1\le C\sum_{j=1}^k\binom{k-1}{j-1}C^j(2b)^k\le C(Cb)^k,\\
 \|\bm a^{2,[k]}\|+\|\bm h^{2,[k]}\|\le C\initscale^2(Cb)^k,\qquad \|\bm z^{[k]}\|\le C\initscale^3(C_0\initscale^4)^k .
\end{gather*}

\emph{(ii), tangents, remainders, absolute sums.} Differentiating along $v$ replaces one factor of each product by its tangent, or one
$(u\psi)^{[a]}$ in a softmax term, or adds $D^{j+1}\omega[\partial_v(u\psi)^{[0]},\ldots]$. With $\|\partial_v\bm h^{0,[0]}\|=O(1)$, $\|\partial_v\psi^{\ell,[0]}\|=O(\initscale^3)$
and $j+1\le2^k$ absorbed into $C_0^k$,
\begin{gather*}
 \|\partial_v\bm\omega^{\ell,[k]}\|_1\le C\initscale^{-1}(C_0\initscale^4)^k,\qquad \|\partial_v\bm a^{\ell,[k]}\|+\|\partial_v\bm h^{\ell,[k]}\|\le C\initscale(C_0\initscale^4)^k,\\
 \|\partial_v\bm z^{[k]}\|\le C\initscale^2(C_0\initscale^4)^k .
\end{gather*}
The integral remainder is bounded by the supremum of the next normalized derivative. At $u=0$,
$\|m^{-1}\sum_j\prod_rF_r(j)\|_{\mathrm{abs}}\le\prod_r\max_j\|F_r(j)\|_{\mathrm{abs}}$, and an order-$k$ monomial has $O(k+1)$ parameter and input
occurrences, so differentiating its norm bound costs at most $\frac{C(k+1)}{\initscale}\le \frac{C2^k}{\initscale}$.

\emph{(iii).} At zero scores the network is affine in the inputs:
\begin{gather*}
 \bm h^1_j=\mathbf M_1\Bigl(\bm x_j+\frac{1}{m_j}\sum_{l\le j}\mathbf W^1_V\bm x_l\Bigr)+\bm o_1,\\
 \bm z^{[0]}_i=\mathbf V\Bigl[\mathbf M_2\Bigl(\bm h^1_i+\frac{1}{m_i}\sum_{j\le i}\mathbf W^2_V\bm h^1_j\Bigr)+\bm o_2\Bigr].
\end{gather*}
Transported terms contain $\mathbf V$, $\mathbf W_V$ and some $\bm x_j$, so their tangent is $O(\initscale^2)$; direct-input and input-free terms have tangent $O(\initscale)$.
With $d(F)$ the number of input occurrences in $F$,
\begin{alignat*}{2}
 d(\bm h^{\ell,[a]})&\le2a+1 &\qquad& \why{induction on $\ell$}\\
 d(\psi^{\ell,[b]})&\le2b+2 && \why{score is a product of two residuals}\\
 d\bigl((u\psi)^{[a]}\bigr)&\le2a && \why{$(u\psi)^{[a]}=\psi^{[a-1]}$}\\
 d(\text{order-$k$ head term})&\le2a_1+\cdots+2a_j+2c+1=2k+1 && \why{$a_1+\cdots+a_j+c=k$}
\end{alignat*}
and affine MLPs and parameter derivatives add none. The one-score correction is
\[
 \frac{1}{m_i}\sum_{j\le i}\bigl(\psi_{ij}(0)-\bar\psi_i(0)\bigr)\mathbf W_V\bm h_j(0),\qquad \bar\psi_i(0)=\frac{1}{m_i}\sum_{l\le i}\psi_{il}(0);
\]
its logit terms contain $\mathbf V,\mathbf W_V,\mathbf W_Q,\mathbf W_K$ and three residual inputs, so they have size $O(\initscale^7)$ and tangent $O(\initscale^6)$.

\emph{(iv).} Fix $l<i$; (iv) concerns the coefficients at $u=0$ and their tangents. Replace $\bm x_l$ by $\chi\bm x_l$. A monomial containing $\bm x_l$
exactly $d\ge1$ times carries $\chi^d$, so
\[
 \sum_{\term\ni\bm x_l}\|\term\|_{\mathrm{abs}}\le\frac{d}{d\chi}\Bigl(\sum_{\term}\|\term\|_{\mathrm{abs}}\Bigr)\Big|_{\chi=1},
\]
a tangent along an input direction, which obeys the product and composition rules above. Call $I^{(\ell)}_j$ a weight of position $j$ if every
marked order-$k$ coefficient of $\bm h^\ell_j$ (absolute sum and tangent) is at most $I^{(\ell)}_j$ times the unmarked bound of (ii); then
$I^{(0)}_j=\mathbf 1[j=l]$. At $u=0$, $\omega_{jj'}=\frac{1}{m_j}$, and each term of a $u$-derivative of $\omega_{jj'}$ is $\omega_{jj'}$ times score entries and
$\E_\omega$-averages, so each entry of an order-$a$ row coefficient is at most $\frac{1}{m_j}$ times the bound \cref{eq:softmax-avg}, which stays
valid with one $\|\bm\nu\|_\infty$ replaced by $\E_\omega|\bm\nu|$. A marked factor of the head output at $j$ contributes:
\[
\begin{array}{@{}lll@{}}
 \text{Marked factor} & \text{Weight} & \text{Reason}\\[2pt]
 \text{value at }j'\le j & \frac{1}{m_j}I^{(\ell-1)}_{j'} & \text{entrywise row bound}\\[2pt]
 \text{key at }j'\le j & \frac{2}{m_j}I^{(\ell-1)}_{j'} & \E_\omega|\bm\nu|=\frac{|\nu_{j'}|}{m_j}\\[2pt]
 \text{query at }j & I^{(\ell-1)}_j & \text{all score entries}
\end{array}
\]
With the residual path and $\|\mathbf M_\ell\|\le2$,
\begin{equation}
 I^{(\ell)}_j\le C\Bigl(I^{(\ell-1)}_j+\frac{1}{m_j}\sum_{j'\le j}I^{(\ell-1)}_{j'}\Bigr).
 \label{eq:weight-rec}
\end{equation}
Unrolling, for $i>l$,
\begin{alignat*}{2}
 I^{(1)}_j&\le C\bigl(\mathbf 1[j=l]+\frac{\mathbf 1[j>l]}{m_j}\bigr) &\qquad& \why{causality, \cref{eq:weight-rec}}\\
 I^{(2)}_i&\le C\Bigl(\frac{C}{m_i}+\frac{1}{m_i}\Bigl(C+\sum_{l<j\le i}\frac{C}{m_j}\Bigr)\Bigr)\le C\kappa_{l,i} && \why{\cref{eq:weight-rec}, \cref{eq:dilution}}
\end{alignat*}
The readout and parameter tangents do not move positions.
\end{proof}


\section{Retaining degreewise credit with conditional projections}\label{sec:retain-credit}
For a function $f$ of the independent letters, let $E_t f$ average over $w_t$ while fixing all other letters, and let $D_t=I-E_t$. Define
\[
 f_S=\left(\prod_{t\in S}D_t\right)\left(\prod_{t\notin S}E_t\right)f,\qquad
 (f)_r=\sum_{|S|=r}f_S.
\]
These operators commute. Expanding $\prod_t(E_t+D_t)=I$ gives $f=\sum_S f_S$. Each $f_S$ depends only on $S$ and has conditional mean zero in each of its letters. If $S\ne U$, choose a letter in their symmetric difference and average it first; this proves $\mathbb E\langle f_S,g_U\rangle=0$. Thus the decomposition is orthogonal. Its uniqueness follows by applying its defining projection to any decomposition with the same conditional-mean properties. This construction is the familiar orthogonal interaction decomposition, retained only for the degree claims.

\begin{lemma}[Matching-degree identity and exact target spectrum]\label{lem:target-contraction}
With $A_r=\mathbb E\|(G)_r\|^2$ and $J_r=\mathbb E\|(J)_r\|_F^2$,
\begin{equation}
 \mathbb E[J^\top G]=\sum_{r=0}^n\mathbb E[(J)_r^\top(G)_r],\qquad
 \|\mathbb E[(J)_r^\top(G)_r]\|\le\sqrt{A_rJ_r}.
 \label{eq:matching}
\end{equation}
The exact target masses are
\[
 \mathbb E\|G_S\|^2=\frac1K\sum_{j=1}^{K-1}|\rho_j|^{2(n-|S|)}(1-|\rho_j|^2)^{|S|},\qquad
 A_r=\binom nr\frac1K\sum_{j=1}^{K-1}|\rho_j|^{2(n-r)}(1-|\rho_j|^2)^r.
\]
In particular $\|G_S\|_{L^2}\le\lambda^{n-|S|}$.
\end{lemma}
\begin{proof}
Orthogonality eliminates all unequal supports and hence unequal degrees, giving the first identity. Cauchy--Schwarz after testing any unit parameter direction gives its bound. For the target, expand
\[
 \alpha^{js_n}=\alpha^{js_0}\prod_t\bigl[\rho_j+(\alpha^{jb_t}-\rho_j)\bigr].
\]
The term indexed by $S$ is $\alpha^{js_0}\rho_j^{n-|S|}\prod_{t\in S}(\alpha^{jb_t}-\rho_j)$. It is centered in every included letter and independent of excluded ones, so it is exactly the projection component. Its squared $L^2$ norm is $|\rho_j|^{2(n-|S|)}(1-|\rho_j|^2)^{|S|}$ by independence. Parseval across the nonconstant Fourier modes gives the vector formula, and summing the $\binom nr$ supports gives $A_r$. The convention is $0^0=1$.
\end{proof}

\begin{proposition}[Score count and degree]
\label{prop:score-degree}
Fix the start state and the query. For $\bm\theta\in\initball$, $\Jz=\sum_{k=0}^{k_0}\mathbf J^{[k]}+\mathbf J_{\mathrm{rem}}$, and there is $\attconst$ such that:
\begin{enumerate}[label=(\roman*),leftmargin=*,itemsep=2pt,topsep=2pt]
\item $\sup\|\mathbf J^{[0]}\|_F\le C\initscale$, $\ \sum_{|S|=1}\sup\|\mathbf J^{[0]}_S\|_F\le C\initscale^2$, and $\mathbf J^{[0]}_S=0$ for $|S|\ge2$;
\item for $k\ge1$, $\ \sum_S\sup\|\mathbf J^{[k]}_S\|_F\le C\initscale^2(\attconst\initscale^4)^k$, and $\mathbf J^{[k]}_S=0$ for $|S|>2k+1$;
\item $\sup\|\mathbf J_{\mathrm{rem}}\|_F\le C\initscale^2(\attconst\initscale^4)^{k_0+1}$.
\end{enumerate}
The same bounds hold for $\bm z^{[k]}$ and $\bm z_{\mathrm{rem}}$ with one more factor $\initscale$. Suprema are over inputs; constants do not depend on
$n$, $k$ or $k_0$.
\end{proposition}

\begin{proof}[Proof of \cref{prop:score-degree}]
For a term $\term$ supported on the sites $U$,
\[
 \term_S=0\quad(S\not\subseteq U),\qquad \sup\|\term_S\|\le2^{|S|}\sup\|\term\|,
\]
so
\[
 \sum_S\sup\|\term_S\|\le\sum_{S\subseteq U}2^{|S|}\sup\|\term\|=3^{|U|}\sup\|\term\| .
\]

\noindent(i) At the answer position the direct input is fixed; apply \cref{prop:network}(iii).

\smallskip\noindent(ii) For $k\ge1$, \cref{prop:network}(ii)--(iii) give $|U|\le2k+1$, hence $\mathbf J^{[k]}_S=0$ for $|S|>2k+1$, and, with $\attconst\ge9C_0$,
\begin{alignat*}{2}
 \sum_S\sup\|\mathbf J^{[k]}_S\|_F&\le\sqrt K\,3^{2k+1}\|\partial\bm z^{[k]}\|_{\mathrm{abs}} &\qquad& \why{support bound}\\
 &\le C\initscale^2(\attconst\initscale^4)^k && \why{\cref{prop:network}(ii)}
\end{alignat*}

\noindent(iii) Apply \cref{prop:network}(ii) with $k=k_0+1$. Logit bounds carry one more factor $\initscale$ than Jacobian bounds.
\end{proof}


\begin{proof}[Complete proof of \cref{thm:main}]
Set $N=\lceil(n-1)/2\rceil$ and $b=C_{\rm att}\varepsilon^4$. Choose $c$ so that $\varepsilon^2\le c\lambda$ implies $b\le\lambda^2/2$. For every support $S$,
\[
 \|\mathbb E[(J^{[k]})_S^\top G_S]\|
 \le\lambda^{n-|S|}\sup\|(J^{[k]})_S\|_F.
\]
The zero-order coefficient is additive, so its projections at degrees at least two vanish. For $k\ge1$, \cref{prop:score-degree} bounds the sum over supports by $C\varepsilon^2b^k\lambda^{n-2k-1}$. Summing over $k$ gives
\[
 C\varepsilon^6\lambda^{n-3}\sum_{k=1}^N(b/\lambda^2)^{k-1}
 \le C\varepsilon^6\lambda^{n-3}.
\]
For the remainder let $Q_{\ge2}$ denote the orthogonal projection onto degrees at least two. Then
\[
 \mathbb E[(Q_{\ge2}J_R)^\top Q_{\ge2}G]
 =\mathbb E[J_R^\top Q_{\ge2}G].
\]
Its norm is at most $\|J_R\|_{L^2(F)}\|Q_{\ge2}G\|_{L^2}\le C\varepsilon^2b^{N+1}$. As in the direct summation proof this is at most $C\varepsilon^6\lambda^{n-3}(b/\lambda^2)^N$. This proves the aggregate interaction bound. Replacing $Q_{\ge2}$ by any single degree projection gives the same bound degree by degree, with the same constant. No estimate of a remainder's supremum after projection is needed.
\end{proof}

\begin{lemma}[Gradient split]\label{lem:gradient-split}
For $p_\theta=\operatorname{softmax}(z_\theta)$,
\[
 \nabla_\theta L_{\rm out}=\mathbb E[J^\top(p_\theta-u)]-\nabla_\theta\outcredit,
 \qquad \nabla_\theta\outcredit=\mathbb E[J^\top G].
\]
\end{lemma}
\begin{proof}
Differentiate the cross-entropy, whose logit derivative is $p_\theta-y$, and split this derivative as $(p_\theta-u)-G$. Finite data support and local smoothness justify differentiation inside the expectation.
\end{proof}

\subsection{The local innovation}
For a uniform start, every $s_{t-1}$ is uniform and independent of the entire letter sequence. Let $q_s=(1-p)e_s+pe_{s+1}$ and
\[
 g_t=e_{s_t}-q_{s_{t-1}}=(b_t-p)(e_{s_{t-1}+1}-e_{s_{t-1}}).
\]
Its mean is zero. In a unitary Fourier basis on $\bm1^\perp$, its covariance has eigenvalues
\[
 \frac{p(1-p)}K|\alpha^j-1|^2
 =\frac{1-|\rho_j|^2}K\ge\frac{1-\lambda^2}K.
\]
Indeed the uniform state gives covariance $I/K$ for its Fourier components, and the independent centered match bit contributes $p(1-p)$. This proves the local innovation inequality without a degree decomposition. Define $\rulecredit=\mathbb E_t\mathbb E\langle z_t,g_t\rangle$.

\subsection{Step rule credit}

Let $\bm g_t=(b_t-p)(\bm e_{s_{t-1}+1}-\bm e_{s_{t-1}})$ and $\rulecredit(\bm\theta)=\E_t\E[\langle\bm z_t,\bm g_t\rangle]$, $t$ uniform on $\{2,\ldots,n\}$. Letter $w_t$ sits at
$l_t=1+t$, and $i_t=n+1+t$ holds $s_{t-1}$ and predicts $s_t$, so $m_{l_t}=t+2$ and $m_{i_t}=n+2+t$. Write $\kappa_t=\kappa_{l_t,i_t}$. The displayed state tokens visible from $i_t$ are functions of $(s_0,w^\star,w_{<t})$; the prompt also includes the letters. The current letter is therefore the only visible source that depends on $b_t$ after conditioning on $s_0,w^\star,w_{\ne t}$.

\begin{lemma}[Average dilution]
\label{lem:dilution-average}
For $n\ge2$,
\begin{equation}
 \E_t[\kappa_t]\le\frac{1}{n+4}\Bigl(1+\log2+\frac{n+2}{n-1}\Bigr)\le\frac{6}{n} .
 \label{eq:dilution-avg}
\end{equation}
\end{lemma}

\begin{proof}
$m_{i_t}\ge n+4$ and $\sum_{l_t<j\le i_t}m_j^{-1}\le\log2+\log\frac{n+2}{t+2}$; also $\sum_{t=2}^n\log\frac{n+2}{t+2}\le\sum_{k=1}^{n+2}\log\frac{n+2}{k}\le n+2$.
\end{proof}

\begin{lemma}[Rule-credit sites]
\label{lem:rule-sites}
Under (A3$'$), a term $\term$ of $\bm z_t$ or $\bm h^2_{i_t}$ has $\E\langle\term,\bm g_t\rangle=0$ unless it contains both the letter input $\bm x_{l_t}$ and the input of a
state ($s_0$ or a displayed state). In particular $\bm z^{[0]}_t$ contributes nothing.
\end{lemma}

\begin{proof}
If $\bm x_{l_t}\notin\term$, then $\E[\bm g_t\mid s_0,w^\star,w_{\ne t}]=0$. If $\term$ contains no state input, it is a function of $(w^\star,w)$ and, by (A3$'$),
$\E[\bm g_t\mid w^\star,w]=(b_t-p)(\bm u-\bm u)=0$. A contributing term has two sites, so $k\ge1$ by \cref{prop:network}(iii).
\end{proof}

For readout indices $a,b$ let
\begin{gather*}
 Q_t=\sum_{\ell=1,2}\sum_{\mathrm{heads}}\frac{1}{m_i}\sum_{j\le i}\mathbf W_V\bm x_j\Bigl(\hat\psi_{ij}-\frac{1}{m_i}\sum_{l\le i}\hat\psi_{il}\Bigr),\qquad i=i_t,\\
 \hat\psi_{ij}=\frac{(\mathbf W_Q\bm x_i)^\top\mathbf W_K\bm x_j}{\sqrt{d_{\mathrm h}}},
\end{gather*}
and $P_n(\tilde{\bm\theta})=\E_t\E[(\bm g_t)_a(Q_t)_b]|_{\bm\theta=\tilde{\bm\theta}}$.

\begin{lemma}[Leading rule coefficient]
\label{lem:leading}
Assume (A1) with full-prefix attention, (A2) and (A3$'$). Fix a readout entry $V_{ab}$ and let
\[
 F_n(\initscale)=\pd{\rulecredit}{V_{ab}}\Big|_{\bm\theta=\initscale\tilde{\bm\theta}}=\E_t\E[(\bm g_t)_ah^2_{i_t,b}]\big|_{\bm\theta=\initscale\tilde{\bm\theta}},\qquad\text{so }\ \|\nabla_{\bm\theta}\rulecredit(\bm\theta_0)\|\ge|F_n(\initscale)|.
\]
\begin{enumerate}[label=(\roman*),leftmargin=*,itemsep=0pt,topsep=2pt]
\item $P_n$ is a homogeneous polynomial of degree $6$ in the attention and embedding entries of $\tilde{\bm\theta}$, independent of the activation
pattern, and for $\initscale\le\initscale_0$ with $\attconst\initscale^4\le\frac{1}{2}$,
\[
 F_n(\initscale)=\initscale^6P_n(\tilde{\bm\theta})+r_n(\initscale),\qquad |r_n(\initscale)|\le\frac{C}{n}\Bigl(\frac{\initscale^7}{\initscale_0}+\initscale^{10}\Bigr).
\]
\item For a fixed query with $0<p<\frac{1}{2}$, $P_n$ contains a degree-$6$ count--letter monomial with coefficient at least
$\dfrac{p(1-p)\min\{1-2p,\frac{1}{4}\}}{3Kn\sqrt{d_{\mathrm h}}}>0$ for $n\ge2$; in particular $P_n\not\equiv0$.
\end{enumerate}
\end{lemma}

\begin{proof}
(i) Let $k_0=2+\lceil\log_2n\rceil$. By \cref{lem:rule-sites}, the zero-score term contributes nothing; by \cref{prop:network} and
\cref{lem:dilution-average}, the terms with at least two scores and the remainder contribute at most
\[
 C\initscale^2\E_t[\kappa_t]\sum_{k=2}^{k_0}(\attconst\initscale^4)^k+C\initscale^2(\attconst\initscale^4)^{k_0+1}\le\frac{C\initscale^{10}}{n} .
\]
With the activation matrices fixed, a one-score monomial satisfies $\term(\initscale\tilde{\bm\theta})=\initscale^{\deg\term}\term(\tilde{\bm\theta})$ with $\deg\term\ge6$, and the
degree-$6$ part is exactly $Q_t$. An MLP matrix correction adds two degrees and a bias output in place of an input adds one, so the
remaining one-score terms are $O(\initscale^7)$; with $\initscale_0\le(2N_b)^{-1}$ and \cref{lem:dilution-average}, their absolute sum is at most
$\frac{C\initscale^7}{n\initscale_0}$, with $C$ independent of $N_b$.

(ii) Choose a block-$2$ head, a state $s^\ast$, an output coordinate $b$, and $a=s^\ast+1$. Set
$\tilde E_{s^\ast,0}=\tilde E_{w^\star,1}=\tilde E_{w^\star,2}=(\tilde{\mathbf W}_Q)_{0,0}=(\tilde{\mathbf W}_K)_{0,1}=(\tilde{\mathbf W}_V)_{b,2}=1$ and all other entries to $0$; the value of $P_n$ is then the
coefficient of the product of these six entries. With $N_t$ the number of visible query tokens,
\begin{gather*}
 (Q_t)_b=\frac{\mathbf 1[s_{t-1}=s^\ast]}{\sqrt{d_{\mathrm h}}}\Bigl(\frac{N_t}{n+2+t}-\frac{N_t^2}{(n+2+t)^2}\Bigr),\\
 \E[(b_t-p)N_t]=p(1-p),\qquad \E[(b_t-p)N_t^2]=p(1-p)\bigl[3+2(n-1)p\bigr].
\end{gather*}
Under (A3$'$), $s_{t-1}$ is independent of the letters, so the coefficient is
\[
 \frac{1}{K\sqrt{d_{\mathrm h}}(n-1)}\sum_{t=2}^n\frac{p(1-p)}{n+2+t}\Bigl(1-\frac{3+2(n-1)p}{n+2+t}\Bigr)
 \ge\frac{p(1-p)\min\{1-2p,\frac{1}{4}\}}{3Kn\sqrt{d_{\mathrm h}}},
\]
using $3+2(n-1)p\le\max\{2p,\frac{3}{4}\}(n+4)$ for $p<\frac{1}{2}$, $n\ge2$, and $n+2+t\le3n$.
\end{proof}

\begin{proof}[Proof of \cref{thm:separation}(i)]
\emph{Step gradient.} At the position $i$ predicting $s_t$, $\bm h^{(0)}_i=\mathbf E_{\mathrm{val}}^\top\bm e_{s_{t-1}}+\mathbf P\bm e_i$. Let
$\bm\Gamma=\frac{1}{K}\sum_{s}\bigl((1-p)\bm e_s+p\bm e_{s+1}-\bm u\bigr)\bm e_s^\top$. For $\initscale\le\frac{\mixrate\embconst}{4KC}$,
\begin{alignat*}{2}
 \bm h_i&=\bm h^{(0)}_i+\bm O^2+O(\initscale^2),\qquad \bm p_i-\bm u=O(\initscale^2) &\qquad& \why{\cref{prop:network}(i)}\\
 \pd{\Lstep}{\mathbf V}&=-\E\bigl[(\bm e_{s_t}-\bm u)\bm h^{(0)\top}_i\bigr]+O(\initscale^2) && \why{$\bm O^2$ fixed, $\E[\bm e_{s_t}-\bm u]=0$}\\
 &=-\bm\Gamma\mathbf E_{\mathrm{val}}+O(\initscale^2) && \why{uniform start}\\
 \|\nabla_{\bm\theta}\Lstep\|&\ge\|\bm\Gamma\mathbf E_{\mathrm{val}}\|_F-C\initscale^2 && \\
 &\ge\frac{\mixrate}{K}\cdot\frac{\initscale\embconst}{2}-C\initscale^2\ge\frac{\initscale\mixrate\embconst}{4K} && \why{$\sigma_{\max}(\bm\Gamma)=\tfrac{\mixrate}{K}$, $\sigma_{\min}(\mathbf E_{\mathrm{val}})\ge\tfrac{\initscale\embconst}{2}$}
\end{alignat*}
using that $\mathbf E_{\mathrm{val}}$ has full row rank.

\emph{Outcome gradient.} On $\initball$ and for $\initscale\le c\mixrate$,
\begin{alignat}{2}
 \nabla_{\bm\theta}\Lout&=\E\bigl[(\Jz)^\top(\bm p-\bm u)\bigr]-\nabla_{\bm\theta}\outcredit &\qquad& \why{\cref{lem:gradient-split}}\nonumber\\
 \|\nabla_{\bm\theta}\Lout\|&\le C\initscale\cdot C\initscale^2+C\initscale\mixrate^n\le C\initscale(\initscale^2+\mixrate^n) && \why{\cref{prop:network}(i), \cref{lem:score-sum}}\label{eq:outcome-grad}
\end{alignat}
\end{proof}

\begin{proof}[Proof of \cref{thm:separation}(ii)]
Let $k_0=1+\lceil\log_2n\rceil$ and $\attconst\initscale^4\le\frac{1}{2}$. By \cref{lem:rule-sites}, only terms with at least one score contribute.
\begin{alignat*}{2}
 \|\nabla_{\bm\theta}\rulecredit(\bm\theta)\|
 &\le2C\initscale^2\,\E_t[\kappa_t]\sum_{k=1}^{k_0}(\attconst\initscale^4)^k+2C\initscale^2(\attconst\initscale^4)^{k_0+1} &\qquad& \why{\cref{prop:network}(ii),(iv)}\\
 &\le C\attconst\initscale^6\bigl(\E_t[\kappa_t]+2^{-k_0}\bigr) && \why{geometric sum}\\
 &\le C\attconst\initscale^6\Bigl(\frac{6}{n}+\frac{1}{2n}\Bigr)\le\frac{C\initscale^6}{n} && \why{\cref{lem:dilution-average}}
\end{alignat*}
\end{proof}

\paragraph{Zero readout.}\label{rem:zero-readout} \Cref{thm:separation} holds with a zero readout: part~(i) uses $\bm p-\bm u=O(\initscale^2)$, exact at $\bm\theta_0$, and part~(ii) concerns $P_n$ of
\cref{lem:leading}, which contains no readout entry.


\begin{remark}[Generic initialization]
\label{rem:generic-rule}
For $p<\frac{1}{2}$, $P_n$ is a nonzero polynomial (\cref{lem:leading}(ii)), so $P_n(\tilde{\bm\theta})\ne0$ almost surely under any absolutely continuous law on
the attention and embedding entries. If $\tilde{\bm\theta}$ satisfies (A2) and $|P_n(\tilde{\bm\theta})|\ge\frac{\rulemargin}{n}$, then for
$\initscale\le\min\{\initscale_0,(2\attconst)^{-1/4},\frac{\rulemargin\initscale_0}{4C},(\frac{\rulemargin}{4C})^{1/4}\}$, with $C$ from \cref{lem:leading}(i),
\[
 \|\nabla_{\bm\theta}\rulecredit(\initscale\tilde{\bm\theta})\|\ge|F_n(\initscale)|\ge\initscale^6\frac{\rulemargin}{n}-\frac{C}{n}\Bigl(\frac{\initscale^7}{\initscale_0}+\initscale^{10}\Bigr)\ge\frac{\initscale^6\rulemargin}{2n},
\]
which matches \cref{thm:separation}(ii) up to constants. Both the threshold and the constant depend on $\rulemargin$; a lower bound on $\rulemargin$
at random initialization would need anti-concentration of $nP_n$, which we do not prove. The manuscript experiments measure the normalized credit.
\end{remark}


\section{Complete binary joint-flow proof}
\label{app:joint}

\providecommand{\why}[1]{\text{\small(#1)}}

\Cref{tab:J-map} lists the lemmas. In displayed chains, the right-hand column names the inequality or lemma used on that line.

\begin{table}[h]
\centering\small
\fitwidth{\begin{tabular}{@{}lll@{}}
\toprule
Lemma & Statement & Uses\\
\midrule
\ref{lem:J-features} & logits $\initscale \mathbf U\bm x_t+\initscale v\,r_{b_t}\mathbf W\bm x_t$ & seed\\
\ref{lem:J-symmetry} & $\mathbf U,\mathbf W$ stay circulant & \ref{lem:J-features}\\
\ref{lem:J-reduce} & autonomous ODE for $(c,g,a)$ & \ref{lem:J-features}, \ref{lem:J-symmetry}\\
\ref{lem:J-geometry} & partial derivatives; convexity & \ref{lem:J-reduce}\\
\ref{lem:J-levels} & loss levels & \ref{lem:J-geometry}\\
\ref{lem:J-gradient} & $\ddt{\mathcal L}\le-\tfrac{1}{3}\initscale^2v^2\tanh^2(\frac{a}{2})(\mathrm{err}_0^2+\mathrm{err}_1^2)$ & \ref{lem:J-geometry}\\
\ref{lem:J-align} & below $H_p$: $ag<0$ and $|a|$ nondecreasing & \ref{lem:J-geometry}, \ref{lem:J-levels}\\
\ref{lem:J-prior} & Phase I: prior fitted by $t_1\le\frac{2h^\star}{\kappa_0\initscale^3}$ & \ref{lem:J-geometry}, \cref{eq:J-eps}\\
\ref{lem:J-phase2} & Phase II: gate aligned and $\mathcal L(t_2)<H_p$ by $t_2$ & \ref{lem:J-geometry}, \ref{lem:J-align}, \ref{lem:J-prior}, \cref{eq:J-eps}\\
\ref{lem:J-descent} & Phase III: margins and $\mathcal L\to0$ & \ref{lem:J-levels}, \ref{lem:J-gradient}, \ref{lem:J-align}\\
\bottomrule
\end{tabular}}
\caption{Lemmas of the proof of \cref{thm:joint}.}
\label{tab:J-map}
\end{table}

\paragraph{Global existence.} The finite-dimensional parameter loss is smooth everywhere. On every interval of existence, $\int_0^t\|\dot\vartheta\|^2\,ds\le\log2$, hence $\|\vartheta(t)-\vartheta(0)\|\le\sqrt{t\log2}$. No finite-time escape is possible, so local uniqueness extends to every finite time. This also justifies the symmetry argument and all stopping-time bootstraps below.

\subsection{Model and exact features}

We count modulo $K$ with a fixed query of probability $p\in(0,1)$ and a uniform start; every non-query letter has the same embedding. Let $\bm x_t=\bm e_{s_{t-1}}$, $\mathbf R\bm e_s=\bm e_{s+1}$, $r_b=\sigma\bigl(-(-1)^ba\bigr)$ for $b\in\{0,1\}$ and $\sigma(z)=(1+e^{-z})^{-1}$.

\paragraph{Seed.} Fix orthonormal residual directions $\hat{\bm e}_0,\hat{\bm e}_\ell$, a count register $(\hat{\bm e}^{\mathrm c}_s)_{s\in\mathbb Z_K}$, a gate register
$(\hat{\bm e}^{\mathrm g}_s)_{s\in\mathbb Z_K}$, and a unit vector $\hat{\bm e}_{\mathrm h}$ of the head space. The embeddings are
\[
 E(s)=\initscale\,\hat{\bm e}^{\mathrm c}_s,\qquad E(\text{query letter})=\initscale(\hat{\bm e}_0-\hat{\bm e}_\ell),\qquad E(\text{other letters})=\initscale(\hat{\bm e}_0+\hat{\bm e}_\ell).
\]
The gate head is head $1$ of block $2$. At the position $i_t$ that holds $s_{t-1}$ and predicts $s_t$ ($t\ge2$), its fixed mask admits
$i_t$ and the position $j_t$ of the current letter, and its relative bias is zero. Its maps, with scalar factors $q,k,v$ and fixed directions, are
\[\mathbf W_Q=q\,\bm\Pi_Q,\qquad \mathbf W_K=k\,\bm\Pi_K,\qquad \mathbf W_V=v\,\bm\Pi_V.\]
The direction matrices are
\[
 \bm\Pi_Q=\textstyle\sum_s\hat{\bm e}_{\mathrm h}(\hat{\bm e}^{\mathrm c}_s)^{\!\top},\qquad \bm\Pi_K=\hat{\bm e}_{\mathrm h}\hat{\bm e}_\ell^{\top},\qquad \bm\Pi_V=\sum_s\hat{\bm e}^{\mathrm g}_s(\hat{\bm e}^{\mathrm c}_s)^{\!\top}.
\]
All other embeddings, heads and MLP weights are zero, and the MLP biases are $\pm\initscale$. The readout is zero except for the
$K\times K$ blocks $\mathbf U$ (count register) and $\mathbf W$ (gate register).

\paragraph{Training.} The trained parameters are $\bm\vartheta=(\mathbf U,\mathbf W,q,k)$, and $\ddt{\bm\vartheta}=-\nabla_{\bm\vartheta}\Lstep(\bm\vartheta)$ for the Euclidean norm on
$\R^{K\times K}\times\R^{K\times K}\times\R\times\R$. The value factor $v$ and all other parameters are fixed. The four trained parameters follow
this one flow at all times; there are no stages and no stopping rule. The initialization is
\[
 \mathbf U(0)=\mathbf W(0)=0,\qquad \tfrac{1}{2}\initscale\le|q(0)|,\ |k(0)|,\ |v|\le2\initscale,
\]
with arbitrary signs.

\begin{lemma}[Exact features]
\label{lem:J-features}
For every $\bm\vartheta$ and $t\ge2$, the logits of the prediction of $s_t$ are
\[
 \bm z_t=\initscale \mathbf U\bm x_t+\initscale v\,r_{b_t}\,\mathbf W\bm x_t,\qquad a=\frac{\initscale^2qk}{\sqrt{d_{\mathrm h}}}.
\]
\end{lemma}

\begin{proof}
Block $1$ and all MLPs add zero, so each residual is its embedding. The query at $i_t$, the keys and the values are
\begin{gather*}
 \text{query}_{i_t}=\initscale q\,\hat{\bm e}_{\mathrm h},\qquad \text{key}_{j_t}=(-1)^{b_t}\initscale k\,\hat{\bm e}_{\mathrm h},\qquad \text{key}_{i_t}=0,\\
 \text{value}_{i_t}=\initscale v\,\hat{\bm e}^{\mathrm g}_{s_{t-1}},\qquad \text{value}_{j_t}=0 .
\end{gather*}
The scores are $(-1)^{b_t}a$ at $j_t$ and $0$ at $i_t$, so $i_t$ gets weight $r_{b_t}$. The residual at $i_t$ becomes
\[
 \initscale\hat{\bm e}^{\mathrm c}_{s_{t-1}}+\initscale v\,r_{b_t}\,\hat{\bm e}^{\mathrm g}_{s_{t-1}},
\]
and the readout gives $\bm z_t$.
\end{proof}

\begin{lemma}[Cyclic symmetry]
\label{lem:J-symmetry}
$RU(t)=\mathbf U(t)\mathbf R$ and $RW(t)=\mathbf W(t)\mathbf R$ for all $t\ge0$.
\end{lemma}

\begin{proof}
Shifting $s_0$ by one preserves the data law and the match bits, and $\ell(\bm z,j+1)=\ell(\mathbf R^{-1}\bm z,j)$ for the cross-entropy. By \cref{lem:J-features},
\[
 \Lstep\circ S=\Lstep,\qquad S(\mathbf U,\mathbf W,q,k)=(\mathbf R^{-1}UR,\ \mathbf R^{-1}WR,\ q,\ k).
\]
$S$ is a linear isometry, so $S\vartheta(t)$ solves the flow, with $S\vartheta(0)=\bm\vartheta(0)$. By uniqueness, $S\vartheta(t)=\bm\vartheta(t)$.
\end{proof}

\subsection{Reduction for \texorpdfstring{$K=2$}{K = 2}}

Let $K=2$, $\ell(z)=\log(1+e^{-z})$, $H_p=-p\log p-(1-p)\log(1-p)$ and $p_{\min}=\min(p,1-p)$.

\begin{lemma}[Reduced loss and flow]
\label{lem:J-reduce}
Write $\mathbf U=u_0\mathbf I+u_1\mathbf R$, $\mathbf W=w_0\mathbf I+w_1\mathbf R$, and
\[
 c=\initscale(u_0-u_1),\qquad g=\initscale v(w_0-w_1),\qquad m_b=(-1)^b\bigl(c+g\,r_b\bigr).
\]
The prediction of $s_t$ has margin $m_{b_t}$ (target logit minus the other logit), so for every $n$
\[
 \Lstep=\mathcal L(c,g,a):=(1-p)\,\ell(m_0)+p\,\ell(m_1).
\]
With $D=q(0)^2-k(0)^2$,
\begin{equation}
 \begin{gathered}
 \ddt{c}=-\initscale^2\pd{\mathcal L}{c},\qquad \ddt{g}=-\initscale^2v^2\pd{\mathcal L}{g},\qquad \ddt{a}=-\frac{\initscale^4S}{d_{\mathrm h}}\pd{\mathcal L}{a},\\
 S:=q^2+k^2=\sqrt{D^2+\frac{4d_{\mathrm h}a^2}{\initscale^4}}.
 \end{gathered}
 \label{eq:J-flow}
\end{equation}
\end{lemma}

\begin{proof}
For $\bm x_t=\bm e_s$, \cref{lem:J-features} gives $z_{t,s}-z_{t,s+1}=c+g\,r_{b_t}$, and the target is $s+b_t$. This gives the margins.

By \cref{lem:J-symmetry} and $\|\mathbf I\|_F^2=\|\mathbf R\|_F^2=2$, $\langle\mathbf I,\mathbf R\rangle_F=0$:
\begin{gather*}
 \nabla_{\mathbf U}\Lstep=\tfrac{1}{2}\Bigl(\pd{\Lstep}{u_0}\mathbf I+\pd{\Lstep}{u_1}\mathbf R\Bigr),\qquad \pd{\Lstep}{u_0}=-\pd{\Lstep}{u_1}=\initscale\, \pd{\mathcal L}{c},\\
 \ddt{c}=\initscale(\ddt{u_0}-\ddt{u_1})=-\initscale^2 \pd{\mathcal L}{c} .
\end{gather*}
The equation for $g$ is identical. For the gate,
\begin{gather*}
 \ddt{q}=-\frac{\initscale^2k}{\sqrt{d_{\mathrm h}}}\pd{\mathcal L}{a},\qquad \ddt{k}=-\frac{\initscale^2q}{\sqrt{d_{\mathrm h}}}\pd{\mathcal L}{a},\qquad \frac{d}{dt}(q^2-k^2)=0,\\
 \ddt{a}=\frac{\initscale^2}{\sqrt{d_{\mathrm h}}}\Bigl(k\ddt{q}+q\ddt{k}\Bigr)=-\frac{\initscale^4S}{d_{\mathrm h}}\pd{\mathcal L}{a},
\end{gather*}
and $S^2=(q^2-k^2)^2+4q^2k^2$.
\end{proof}

\paragraph{Notation.} $r_0+r_1=1$, $r_1-r_0=\tanh(\frac{a}{2})$ and $\frac{d}{da}\tanh(\frac{a}{2})=2\sigma'(a)$. Let
\begin{gather*}
 \mathrm{err}_0=(1-p)\,\sigma(-m_0),\qquad \mathrm{err}_1=p\,\sigma(-m_1),\qquad \mathrm{err}=\mathrm{err}_0+\mathrm{err}_1,\\
 h=\tfrac{1}{2}(m_0-m_1)=c+\tfrac{g}{2},\qquad \Delta=m_0+m_1=-g\tanh(\frac{a}{2}),
\end{gather*}
the weighted errors, the count-only margin, and the common margin shift from the learned rule; $m_0=h+\frac{\Delta}{2}$ and
$m_1=-h+\frac{\Delta}{2}$.

\begin{lemma}[Derivatives]
\label{lem:J-geometry}
Let $\tilde{\mathcal L}( h,\Delta)=(1-p)\ell( h+\tfrac{\Delta}{2})+p\,\ell(- h+\tfrac{\Delta}{2})$ and $ \mathcal L_0( h)=\tilde{\mathcal L}( h,0)$.
\begin{enumerate}[label=(\alph*),leftmargin=*,itemsep=2pt,topsep=3pt]
\item $\mathcal L=\tilde{\mathcal L}( h,\Delta)$, $\tilde{\mathcal L}$ is jointly convex, and $\pd{\tilde{\mathcal L}}{\Delta}=-\frac{\mathrm{err}}{2}<0$.
\item $ \pd{\mathcal L}{c}=\mathrm{err}_1-\mathrm{err}_0$, $\quad \pd{\mathcal L}{g}=r_1\mathrm{err}_1-r_0\mathrm{err}_0=\tfrac{1}{2}\bigl( \pd{\mathcal L}{c}+\tanh(\frac{a}{2})\,\mathrm{err}\bigr)$, $\quad \pd{\mathcal L}{a}=g\,\sigma'(a)\,\mathrm{err}$.
\item $\kappa:=\frac{\partial^2\tilde{\mathcal L}}{\partial h^2}=(1-p)\sigma'(m_0)+p\sigma'(m_1)$, $\quad|\frac{\partial^2\tilde{\mathcal L}}{\partial h\,\partial\Delta}|\le\min(\tfrac{1}{8},\tfrac{\kappa}{2})$, $\quad| \pd{\mathcal L}{c}- \mathcal L_0'( h)|\le\tfrac{1}{8}|\Delta|$.
\item $ \mathcal L_0'(z)=\sigma(z)-(1-p)$ and $ \mathcal L_0''=\sigma'\le\tfrac{1}{4}$.
\item $\ddt{\mathcal L}=-\initscale^2( \pd{\mathcal L}{c})^2-\initscale^2v^2( \pd{\mathcal L}{g})^2-\frac{\initscale^4S\,( \pd{\mathcal L}{a})^2}{d_{\mathrm h}}\le0$.
\end{enumerate}
\end{lemma}

\begin{proof}
(a) $m_0=h+\frac{\Delta}{2}$, $m_1=-h+\frac{\Delta}{2}$, and $\ell$ is convex and decreasing.

(b) With $\ell'(z)=-\sigma(-z)$: $\pd{\tilde{\mathcal L}}{h}=\mathrm{err}_1-\mathrm{err}_0$ and $\pd{\tilde{\mathcal L}}{\Delta}=-\frac{\mathrm{err}}{2}$. Then use
\begin{gather*}
 \pd{h}{c}=1,\qquad \pd{h}{g}=\tfrac{1}{2},\qquad \pd{\Delta}{g}=-\tanh(\frac{a}{2}),\qquad \pd{\Delta}{a}=-2g\,\sigma'(a),\\
 r_b=\tfrac{1}{2}\bigl(1-(-1)^b\tanh(\frac{a}{2})\bigr).
\end{gather*}

(c) $\frac{\partial^2\tilde{\mathcal L}}{\partial h\,\partial\Delta}=\frac{1}{2}\bigl[(1-p)\sigma'(m_0)-p\sigma'(m_1)\bigr]$, and $\sigma'\le\frac{1}{4}$. The last bound is the mean value theorem in $\Delta$.

(d) is direct. (e) is $\ddt{\mathcal L}= \pd{\mathcal L}{c}\,\ddt{c}+ \pd{\mathcal L}{g}\,\ddt{g}+ \pd{\mathcal L}{a}\,\ddt{a}$ with \cref{eq:J-flow}.
\end{proof}

\begin{lemma}[Loss levels]
\label{lem:J-levels}
Let $\mathcal L^\star=\frac{1}{2}p_{\min}\log2$.
\begin{enumerate}[label=(\alph*),leftmargin=*,itemsep=2pt,topsep=3pt]
\item If $\Delta\le0$, then $\mathcal L\ge H_p$.
\item If $\mathcal L<p_{\min}\log2$, then $m_0>0$ and $m_1>0$: every prediction of $s_t$, $t\ge2$, has positive margin.
\item If $m_0,m_1\ge0$, then $\mathrm{err}_0^2+\mathrm{err}_1^2\ge \frac{\mathcal L^2}{8}$.
\item If $\mathcal L>\mathcal L^\star$, then $\mathrm{err}_0^2+\mathrm{err}_1^2\ge \frac{p_{\min}^2}{70}$.
\end{enumerate}
\end{lemma}

\begin{proof}
(a) $\mathcal L=\tilde{\mathcal L}( h,\Delta)\ge \mathcal L_0( h)\ge\min \mathcal L_0=H_p$, attained at $\sigma( h)=1-p$.

(b) $m_0\le0\implies \mathcal L\ge(1-p)\log2$, and $m_1\le0\implies \mathcal L\ge p\log2$.

(c) For $z\ge0$, $\ell(z)\le e^{-z}\le2\sigma(-z)$. Hence $\mathcal L\le2\mathrm{err}$ and $\mathrm{err}_0^2+\mathrm{err}_1^2\ge\frac{1}{2}\mathrm{err}^2\ge\frac{1}{8}\mathcal L^2$.

(d) If $m_0<0$, then $\mathrm{err}_0>\frac{1}{2}(1-p)$; if $m_1<0$, then $\mathrm{err}_1>\frac{1}{2}p$; otherwise (c) applies. So
\[
 \mathrm{err}_0^2+\mathrm{err}_1^2\ge\min\Bigl\{\frac{p_{\min}^2}{4},\ \frac{(\mathcal L^\star)^2}{8}\Bigr\}=\frac{(\log2)^2}{32}\,p_{\min}^2\ge\frac{p_{\min}^2}{70}.
\]
\end{proof}

\begin{lemma}[Gradient lower bound]
\label{lem:J-gradient}
$( \pd{\mathcal L}{c})^2+( \pd{\mathcal L}{g})^2\ge\frac{1}{3}\tanh^2(\frac{a}{2})\,(\mathrm{err}_0^2+\mathrm{err}_1^2)$. If $\initscale^2v^2\le\initscale^2$, then
\[
 \ddt{\mathcal L}\le-\tfrac{1}{3}\initscale^2v^2\tanh^2(\frac{a}{2})\,(\mathrm{err}_0^2+\mathrm{err}_1^2).
\]
\end{lemma}

\begin{proof}
By \cref{lem:J-geometry}(b), $(\pd{\mathcal L}{c},\pd{\mathcal L}{g})^\top=\mathbf A(\mathrm{err}_0,\mathrm{err}_1)^\top$ with $\mathbf A=\bigl(\begin{smallmatrix}-1&1\\-r_0&r_1\end{smallmatrix}\bigr)$. Its singular values $s_1\ge s_2$
satisfy
\begin{gather*}
 s_1s_2=|\det\mathbf A|=|\tanh(\frac{a}{2})|,\qquad s_1\le\|\mathbf A\|_F=(2+r_0^2+r_1^2)^{1/2}\le\sqrt3,\\
 s_2\ge\frac{|\tanh(\frac{a}{2})|}{\sqrt3} .
\end{gather*}
The second claim follows from \cref{lem:J-geometry}(e).
\end{proof}

\begin{lemma}[Alignment below the prior]
\label{lem:J-align}
\begin{enumerate}[label=(\alph*),leftmargin=*,itemsep=2pt,topsep=3pt]
\item $\operatorname{sign}\ddt{a}=-\operatorname{sign}g$.
\item If $\mathcal L(t_0)<H_p$, then for all $t\ge t_0$: $\Delta(t)>0$, $a(t)g(t)<0$ and $|a(t)|\ge|a(t_0)|$.
\end{enumerate}
\end{lemma}

\begin{proof}
(a) By \cref{eq:J-flow},
\[
 \ddt{a}=-\frac{\initscale^4S\,g\,\sigma'(a)\,\mathrm{err}}{d_{\mathrm h}},\qquad S,\ \sigma'(a),\ \mathrm{err}>0.
\]

(b) $\mathcal L(t)\le \mathcal L(t_0)<H_p$, so $\Delta(t)>0$ by \cref{lem:J-levels}(a). Since $\Delta=-g\tanh(\frac{a}{2})$, $ag<0$, and
\[
 \frac{d}{dt}a^2=-\frac{2\initscale^4S}{d_{\mathrm h}}\,ag\,\sigma'(a)\,\mathrm{err}>0 .
\]
\end{proof}

\paragraph{Reflection.} The map $(c,g,a)\mapsto(-c,-g,-a)$ exchanges $m_0$ and $m_1$, and so maps $\mathcal L$ and \cref{eq:J-flow} for $p$ to those for
$1-p$. On parameters it is $(u_0,u_1,w_0,w_1,q,k)\mapsto(u_1,u_0,w_1,w_0,-q,k)$, i.e.\ $q(0)\mapsto-q(0)$. It preserves $q(0)k(0)(1-2p)$, $D=q(0)^2-k(0)^2$
and the conclusion of \cref{thm:joint}. Hence we assume $p<\frac{1}{2}$, and set
\[
  h^\star=\log\frac{1-p}{p}>0,\qquad \kappa_0=\frac{p(1-p)}{e}\le\frac{1}{4e},\qquad  \mathcal L_0'( h^\star)=0 .
\]
The initialization is \emph{aligned} if $q(0)k(0)<0$, i.e.\ $a(0)<0$.

\subsection{Conditions and scales}
\label{app:J-constants}

\paragraph{Threshold.} Let $\gamma=\mathrm{gap}$ if $q(0)k(0)>0$ and $\gamma=1$ if $q(0)k(0)<0$, and
\[
 \initscale^\star=\frac{\kappa_0\,p(1-p)\,\min(h^\star,1)\,\gamma}{1.3\cdot10^7\,(1+h^\star)^2\sqrt{d_{\mathrm h}}} ;
\]
in the misaligned case also $|D|\ge\mathrm{gap}\,\initscale^2$. Let $0<\initscale\le\initscale^\star$. Since $\initscale^\star<1$, $\min(h^\star,1)\le h^\star$,
$\frac{h^\star}{(1+h^\star)^2}\le\frac{1}{4}$, $\frac{(h^\star)^2}{(1+h^\star)^2}\le1$, $\gamma\le1$, $\initscale^3,\initscale^6\le\initscale$ and $\initscale^2\le(\initscale^\star)^2$,
\begin{equation}
 \begin{gathered}
 \initscale\le\frac{h^\star}{4},\qquad \initscale\le\frac{\kappa_0}{4}\le\frac{1}{40},\qquad \initscale h^\star\le\frac{\kappa_0}{45}\le\frac{1}{8},\qquad
 \initscale(h^\star)^2\le\frac{\kappa_0}{20},\\
 \initscale\le\frac{\kappa_0h^\star}{128},\qquad \initscale^3\le\frac{\kappa_0h^\star}{30720},\qquad
 \initscale^6\le\frac{\kappa_0p(1-p)h^\star}{7.2\cdot10^6(1+h^\star)\sqrt{d_{\mathrm h}}},\\
 \initscale^2\le\frac{\kappa_0\,\mathrm{gap}\,(h^\star)^2}{1.3\cdot10^7\sqrt{d_{\mathrm h}}}\quad\text{if }q(0)k(0)>0,
 \end{gathered}
 \label{eq:J-eps}
\end{equation}
and $\initscale\log\frac{1}{\initscale}\le e^{-1}$. These are the only smallness properties of $\initscale$ used below.

\paragraph{Scales.} Let
\begin{gather*}
 a_{\mathrm{sc}}=\frac{\initscale^4}{\sqrt{d_{\mathrm h}}}\le\initscale^4,\qquad
 T_{\mathrm{relax}}=\frac{5\log\frac{1}{\initscale}}{\kappa_0\initscale^2},\\
 T_{\mathrm{align}}=\begin{cases}\dfrac{2048\sqrt{d_{\mathrm h}}}{\kappa_0h^\star\,\mathrm{gap}\,\initscale^4}&\text{if }q(0)k(0)>0,\\0&\text{if }q(0)k(0)<0.\end{cases}
\end{gather*}
From the seed,
\begin{equation}
 \begin{gathered}
 \tfrac{1}{4}\initscale^4\le \initscale^2v^2\le4\initscale^4\le\initscale^2,\qquad |D|\le4\initscale^2,\qquad \tfrac{1}{4}a_{\mathrm{sc}}\le|a(0)|\le4a_{\mathrm{sc}},\\
 |a|\le ma_{\mathrm{sc}}\implies S\le(4+2m)\initscale^2,\qquad |a|\le1\implies\sigma'(a)\ge\tfrac{1}{6} .
 \end{gathered}
 \label{eq:J-scales}
\end{equation}


\subsection{Phase I: the prior}

\begin{lemma}[Prior fitting]
\label{lem:J-prior}
Let $t_1=\inf\{t\ge0: h(t)\ge h^\star-\initscale\}$. Then $t_1\le\frac{2 h^\star}{\kappa_0\initscale^3}$; on $[0,t_1]$, $c,g, h$ are nondecreasing and $a$ is
nonincreasing; and, with $ g_1=g(t_1)$,
\begin{gather*}
  h(t_1)= h^\star-\initscale,\qquad \tfrac{1}{32}\initscale^2 h^\star\le g_1\le3\initscale^2 h^\star,\\
 |a(t_1)|\le5a_{\mathrm{sc}},\qquad a(t_1)\le a(0),\qquad \Bigl|\pd{\mathcal L}{c}(t_1)\Bigr|\le\tfrac{\initscale}{2} .
\end{gather*}
\end{lemma}

\begin{proof}
Let $[0,t')$ be the maximal subinterval of $[0,t_1)$ on which
\[
 \mathsf B_1:\qquad |a|\le6a_{\mathrm{sc}},\qquad 0\le g\le4\initscale^2 h^\star,\qquad 0\le h< h^\star-\initscale .
\]
$\mathsf B_1$ holds at $t=0$ by \cref{eq:J-scales} and $\initscale\le\frac{h^\star}{4}$. On $[0,t')$,
\begin{alignat*}{2}
 &S\le16\initscale^2,\qquad |\tanh(\frac{a}{2})|\le3\initscale^4,\qquad |\Delta|\le12\initscale^6 h^\star, &\qquad& \\
 &\mathcal L_0'(h)\le-\kappa_0\initscale && \why{$\mathcal L_0''=\sigma'\ge p(1-p)$ on $[0,h^\star]$}
\end{alignat*}
Hence
\begin{align}
  \pd{\mathcal L}{c}&\le \mathcal L_0'( h)+\tfrac{1}{8}|\Delta|\le-\kappa_0\initscale+\tfrac{3}{2}\initscale^6 h^\star\le-\tfrac{1}{2}\kappa_0\initscale && \why{\cref{eq:J-eps}}\label{eq:J-I-Fx}\\
 |\tanh(\frac{a}{2})|\,\mathrm{err}&\le3\initscale^4\le\tfrac{1}{8}\kappa_0\initscale\le\tfrac{1}{4}\Bigl|\pd{\mathcal L}{c}\Bigr| && \why{$\initscale\le\tfrac{\kappa_0}{4}$}\label{eq:J-I-tau}
\end{align}
and
\begin{alignat}{2}
 \ddt{c}&=\initscale^2\Bigl|\pd{\mathcal L}{c}\Bigr| &\qquad& \why{\cref{eq:J-flow}, \cref{eq:J-I-Fx}}\nonumber\\
 \tfrac{3}{8}\initscale^2v^2\Bigl|\pd{\mathcal L}{c}\Bigr|\le\ddt{g}&\le\tfrac{5}{8}\initscale^2v^2\Bigl|\pd{\mathcal L}{c}\Bigr| && \why{\cref{lem:J-geometry}(b), \cref{eq:J-I-tau}}\label{eq:J-I-rates}\\
 \ddt{h}=\ddt{c}+\tfrac{1}{2}\ddt{g}&\ge\initscale^2\Bigl|\pd{\mathcal L}{c}\Bigr|\ge\tfrac{1}{2}\kappa_0\initscale^3 && \why{\cref{eq:J-I-Fx}}\nonumber
\end{alignat}
So $c,g, h$ are nondecreasing and $t'\le\frac{2 h^\star}{\kappa_0\initscale^3}$. By \cref{eq:J-I-rates,eq:J-scales}, $\frac{\ddt{g}}{\ddt{c}}\in[\frac{3}{32}\initscale^2,\frac{5}{2}\initscale^2]$, hence
\begin{equation}
 \tfrac{3}{32}\initscale^2c\le g\le\tfrac{5}{2}\initscale^2c<\tfrac{5}{2}\initscale^2 h^\star .
 \label{eq:J-I-g}
\end{equation}
By \cref{lem:J-align}(a), $\ddt{a}\le0$, and
\begin{align*}
 |\ddt{a}|&\le\frac{\initscale^4S}{d_{\mathrm h}}\cdot\frac{g}{4}\le\frac{10\,\initscale^8 h^\star}{d_{\mathrm h}},\\
 |a(t)-a(0)|&\le\frac{10\,\initscale^8 h^\star}{d_{\mathrm h}}\cdot\frac{2 h^\star}{\kappa_0\initscale^3}\le a_{\mathrm{sc}} && \why{$\initscale(h^\star)^2\le\tfrac{\kappa_0}{20}$}
\end{align*}
With $|a(0)|\le4a_{\mathrm{sc}}$ and $a$ nonincreasing, $|a|\le5a_{\mathrm{sc}}$. So $\mathsf B_1$ holds strictly on $[0,t')$, and $t'=t_1$ by continuity.

At $t_1$, by $\initscale\le\frac{h^\star}{4}$ and \cref{eq:J-I-g},
\[
  h= h^\star-\initscale\ge\tfrac{3}{4} h^\star,\qquad c= h-\tfrac{g}{2}\ge\tfrac{3}{4} h^\star-\tfrac{5}{4}\initscale^2 h^\star\ge\tfrac{1}{2} h^\star,\qquad  g_1\ge\tfrac{3}{64}\initscale^2 h^\star .
\]
Finally $ \mathcal L_0'( h^\star-\initscale)\in[-\frac{\initscale}{4},0]$ by \cref{lem:J-geometry}(d), and $\frac{1}{8}|\Delta|\le\frac{3}{2}\initscale^6 h^\star\le\frac{\initscale}{4}$.
\end{proof}

\subsection{Phase II: alignment of the gate}

Let
\begin{gather*}
 t_{\mathrm{align}}=\inf\Bigl\{t\ge t_1:a(t)\le-\frac{a_{\mathrm{sc}}}{8}\Bigr\},\qquad t_2=\max\{t_1+T_{\mathrm{relax}},\,t_{\mathrm{align}}\},\qquad
 \Delta_2=\frac{\initscale^6h^\star}{2048\sqrt{d_{\mathrm h}}}.
\end{gather*}

\begin{lemma}[Gate alignment]
\label{lem:J-phase2}
$t_2\le t_1+T_{\mathrm{relax}}+T_{\mathrm{align}}$, and on $[t_1,t_2]$
\[
 \Bigl|\pd{\mathcal L}{c}(t)\Bigr|\le22\initscale^6+\frac{\initscale}{2}e^{-\kappa_0\initscale^2(t-t_1)},\qquad |h-h^\star|\le\frac{3\initscale}{\kappa_0},\qquad
 |g-g_1|\le\frac{1}{2}g_1,\qquad |a|\le9a_{\mathrm{sc}} .
\]
At $t_2$,
\[
 \Bigl|\pd{\mathcal L}{c}(t_2)\Bigr|\le23\initscale^6,\qquad a(t_2)\le-\frac{a_{\mathrm{sc}}}{8},\qquad \mathcal L(t_2)\le H_p-\frac{1}{4}p(1-p)\Delta_2<H_p .
\]
\end{lemma}

\begin{proof}
Let $M_\infty=22\initscale^6$ and $M(t)=M_\infty+\frac{\initscale}{2}e^{-\kappa_0\initscale^2(t-t_1)}\le\initscale$. Let $[t_1,t^\ast)$ be the maximal subinterval of
$[t_1,t_2)$ on which
\[
 \mathsf B_2:\qquad \Bigl|\pd{\mathcal L}{c}\Bigr|\le2M,\qquad |h-h^\star|\le1,\qquad |g-g_1|\le\frac{3}{4}g_1,\qquad |a|\le20a_{\mathrm{sc}} .
\]
$\mathsf B_2$ holds strictly at $t_1$ by \cref{lem:J-prior}. Parts (a)--(f) hold on $[t_1,t^\ast)$; (g) closes the bootstrap and (h) proves the loss bound.

\smallskip\noindent(a) \emph{Consequences of $\mathsf B_2$.}
\begin{alignat*}{2}
 \tfrac{1}{4}g_1\le g&\le\tfrac{7}{4}g_1\le\tfrac{21}{4}\initscale^2h^\star\le6\initscale^2h^\star &\qquad& \why{$\mathsf B_2$, \cref{lem:J-prior}}\\
 |\tanh(\tfrac{a}{2})|&\le\tfrac{1}{2}|a|\le10a_{\mathrm{sc}} && \why{$\mathsf B_2$}\\
 |\Delta|&=g\,|\tanh(\tfrac{a}{2})|\le60\initscale^6h^\star\le\initscale && \why{$\initscale h^\star\le\tfrac{1}{8}$, $\initscale\le\tfrac{1}{40}$}\\
 S&\le44\initscale^2,\qquad \Bigl|\pd{\mathcal L}{c}\Bigr|\le2M\le2\initscale,\qquad \sigma'(a)\ge\tfrac{1}{6} && \why{\cref{eq:J-scales}, $|a|\le20\initscale^4\le1$}
\end{alignat*}

\noindent(b) \emph{Location.} Since $\sigma'$ decreases in $|x|$ and $p+e(1-p)\le e$, $\sigma'(x)\ge\sigma'(h^\star+1)=\frac{e\,p(1-p)}{(p+e(1-p))^2}\ge\kappa_0$
for $|x-h^\star|\le1$. With $\mathcal L_0'(h^\star)=0$ and $|h-h^\star|\le1$,
\begin{alignat*}{2}
 \kappa_0|h-h^\star|&\le|\mathcal L_0'(h)|\le\Bigl|\pd{\mathcal L}{c}\Bigr|+\tfrac{1}{8}|\Delta| &\qquad& \why{mean value, \cref{lem:J-geometry}(c)}\\
 &\le3\initscale\le\tfrac{3}{4}\kappa_0 && \why{(a), $\initscale\le\tfrac{\kappa_0}{4}$}
\end{alignat*}
So $m_0=h+\frac{1}{2}\Delta$ and $-m_1=h-\frac{1}{2}\Delta$ lie in $[h^\star-1,h^\star+1]$, and
\begin{alignat*}{2}
 \kappa&=(1-p)\sigma'(m_0)+p\,\sigma'(m_1)\ge\kappa_0 &\qquad& \why{$\sigma'$ even}\\
 \mathrm{err}&\ge(1-p)\,\sigma(-h^\star-1)+p\,\sigma(h^\star-1)\ge(1-p)\frac{p}{e}+p\,\frac{1-p}{e}=2\kappa_0 && \why{monotonicity of $\sigma$}
\end{alignat*}

\noindent(c) \emph{Tracking.} For almost every $t$,
\begin{equation}
 \frac{d}{dt}\Bigl|\pd{\mathcal L}{c}\Bigr|\le-\kappa\initscale^2\Bigl(\Bigl|\pd{\mathcal L}{c}\Bigr|-M_\infty\Bigr).
 \label{eq:J-track}
\end{equation}
Indeed, by \cref{eq:J-flow} and \cref{lem:J-geometry}(b),
\begin{align*}
 \ddt{h}&=-\bigl(\initscale^2+\tfrac{\initscale^2v^2}{4}\bigr) \pd{\mathcal L}{c}-\tfrac{\initscale^2v^2}{4}\tanh(\tfrac{a}{2})\,\mathrm{err},\\
 \ddt{\Delta}&=\tfrac{\initscale^2v^2}{2}\tanh(\tfrac{a}{2})\Bigl( \pd{\mathcal L}{c}+\tanh(\tfrac{a}{2})\mathrm{err}\Bigr)+\frac{2\initscale^4S}{d_{\mathrm h}}\,g^2\sigma'(a)^2\mathrm{err},
\end{align*}
and $\pd{\mathcal L}{c}=\pd{\tilde{\mathcal L}}{h}(h,\Delta)$ gives $\frac{d}{dt}\pd{\mathcal L}{c}=\kappa\,\ddt{h}+\frac{\partial^2\tilde{\mathcal L}}{\partial h\,\partial\Delta}\,\ddt{\Delta}$. With
$|\frac{\partial^2\tilde{\mathcal L}}{\partial h\,\partial\Delta}|\le\frac{\kappa}{2}$, $|\pd{\mathcal L}{c}|+|\tanh(\frac{a}{2})|\le1$, $\mathrm{err}\le1$ and $\sigma'\le\frac{1}{4}$,
\begin{alignat*}{2}
 \frac{d}{dt}\Bigl|\pd{\mathcal L}{c}\Bigr|&\le-\kappa\initscale^2\Bigl|\pd{\mathcal L}{c}\Bigr|+\kappa\Bigl(\frac{\initscale^2v^2}{2}\Bigl|\tanh(\tfrac{a}{2})\Bigr|+\frac{\initscale^4S}{16\,d_{\mathrm h}}\,g^2\Bigr) &\qquad& \\
 \tfrac{\initscale^2v^2}{2}|\tanh(\tfrac{a}{2})|&\le2\initscale^4\cdot10\initscale^4=20\initscale^8 && \why{\cref{eq:J-scales}, (a)}\\
 \tfrac{\initscale^4S}{16d_{\mathrm h}}\,g^2&\le\tfrac{44}{16}\initscale^6\cdot36\initscale^4(h^\star)^2\le2\initscale^8 && \why{(a), $\initscale h^\star\le\tfrac{1}{8}$}
\end{alignat*}
which is \cref{eq:J-track}. With $u=(|\pd{\mathcal L}{c}|-M_\infty)_+$, \cref{eq:J-track} and (b) give $\ddt{u}\le-\kappa_0\initscale^2u$, so
$u(t)\le u(t_1)e^{-\kappa_0\initscale^2(t-t_1)}\le\frac{\initscale}{2}e^{-\kappa_0\initscale^2(t-t_1)}$ by \cref{lem:J-prior}. Hence $|\pd{\mathcal L}{c}|\le M$, and
$|\pd{\mathcal L}{c}|\le23\initscale^6$ once $t-t_1\ge T_{\mathrm{relax}}$, where the exponential is at most $\initscale^5$.

\smallskip\noindent(d) \emph{Elapsed time: $t-t_1\le T_{\mathrm{relax}}+T_{\mathrm{align}}$.} If $q(0)k(0)<0$, then $a(t_1)\le a(0)\le-\frac{1}{4}a_{\mathrm{sc}}$, so
$t_{\mathrm{align}}=t_1$ and $t<t_2=t_1+T_{\mathrm{relax}}$. If $q(0)k(0)>0$ and $t>t_1+T_{\mathrm{relax}}$, then $t<t_{\mathrm{align}}$, so $a>-\frac{a_{\mathrm{sc}}}{8}$ on
$[t_1,t]$, and there
\begin{alignat*}{2}
 S&\ge|D|\ge\initscale^2\mathrm{gap},\qquad g\ge\tfrac{1}{4}g_1\ge\tfrac{1}{128}\initscale^2h^\star &\qquad& \why{$q^2-k^2=D$, (a), \cref{lem:J-prior}}\\
 -\ddt{a}&=\frac{\initscale^4S}{d_{\mathrm h}}\,g\,\sigma'(a)\,\mathrm{err}\ge\frac{\initscale^8\mathrm{gap}\,h^\star}{128\,d_{\mathrm h}}\cdot\frac{1}{6}\cdot2\kappa_0=\frac{\initscale^8\mathrm{gap}\,h^\star\kappa_0}{384\,d_{\mathrm h}} && \why{\cref{lem:J-align}(a), (a), (b)}\\
 t-t_1&\le\frac{41}{8}a_{\mathrm{sc}}\cdot\frac{384\,d_{\mathrm h}}{\initscale^8\mathrm{gap}\,h^\star\kappa_0}=\frac{1968\sqrt{d_{\mathrm h}}}{\kappa_0h^\star\mathrm{gap}\,\initscale^4}\le T_{\mathrm{align}} && \why{$a(t_1)\le5a_{\mathrm{sc}}$}
\end{alignat*}

\noindent(e) \emph{Gate contrast: $|g-g_1|\le\frac{1}{2}g_1$.}
\begin{alignat*}{2}
 |\ddt{g}|&\le\tfrac{\initscale^2v^2}{2}\Bigl(\Bigl|\pd{\mathcal L}{c}\Bigr|+|\tanh(\tfrac{a}{2})|\Bigr)\le2\initscale^4\Bigl(\Bigl|\pd{\mathcal L}{c}\Bigr|+10\initscale^4\Bigr) &\qquad& \why{\cref{eq:J-flow}, (a)}\\
 \int_{t_1}^t\Bigl|\pd{\mathcal L}{c}\Bigr|&\le M_\infty(t-t_1)+\frac{1}{2\kappa_0\initscale} && \why{(c)}\\
 |g-g_1|&\le24\initscale^8(T_{\mathrm{relax}}+T_{\mathrm{align}})+\frac{\initscale^3}{\kappa_0}\le\frac{\initscale^2h^\star}{128}+\frac{\initscale^2h^\star}{128} && \why{(d), \cref{eq:J-eps}}
\end{alignat*}
using $T_{\mathrm{relax}},T_{\mathrm{align}}\le\frac{h^\star}{6144\initscale^6}$ and $\frac{\initscale^3}{\kappa_0}\le\frac{\initscale^2h^\star}{128}$, both from \cref{eq:J-eps} and $\initscale^4\log\frac{1}{\initscale}\le\initscale^3$. So
$|g-g_1|\le\frac{1}{64}\initscale^2h^\star\le\frac{1}{2}g_1$.

\smallskip\noindent(f) \emph{Gate strength: $a$ is nonincreasing and $|a|\le9a_{\mathrm{sc}}$.} By (e), $g>0$, so $\ddt{a}<0$ and $a\le a(t_1)\le5a_{\mathrm{sc}}$. For $t\le t_{\mathrm{align}}$,
$a(t)\ge\min\{a(t_1),-\frac{a_{\mathrm{sc}}}{8}\}\ge-5a_{\mathrm{sc}}$. For $t>t_{\mathrm{align}}$, $t<t_2$ forces $t-t_{\mathrm{align}}\le T_{\mathrm{relax}}$, and
\begin{alignat*}{2}
 |\ddt{a}|&\le\frac{\initscale^4S}{d_{\mathrm h}}\cdot\frac{g}{4}\le\frac{66\,\initscale^8h^\star}{d_{\mathrm h}} &\qquad& \why{(a)}\\
 |a(t)-a(t_{\mathrm{align}})|&\le\frac{330\,(\initscale h^\star)\,(\initscale\log\frac{1}{\initscale})}{\kappa_0}\,a_{\mathrm{sc}}\le\frac{330}{45e}\,a_{\mathrm{sc}}\le4a_{\mathrm{sc}} && \why{$\initscale h^\star\le\tfrac{\kappa_0}{45}$}
\end{alignat*}
so $a(t)\ge-9a_{\mathrm{sc}}$.

\smallskip\noindent(g) \emph{Closing.} By (b), (c), (e), (f), on $[t_1,t^\ast)$
\[
 \Bigl|\pd{\mathcal L}{c}\Bigr|\le M<2M,\qquad |h-h^\star|\le\frac{3}{4}<1,\qquad |g-g_1|\le\frac{1}{2}g_1<\frac{3}{4}g_1,\qquad |a|\le9a_{\mathrm{sc}}<20a_{\mathrm{sc}},
\]
each strictly inside $\mathsf B_2$. If $t^\ast<t_2$, continuity would extend $\mathsf B_2$ past $t^\ast$; so $t^\ast=t_2$, and (d) gives
$t_2-t_1\le T_{\mathrm{relax}}+T_{\mathrm{align}}$. The bounds extend to $t_2$ by continuity, and $a(t_2)\le a(t_{\mathrm{align}})\le-\frac{a_{\mathrm{sc}}}{8}$ since $a$ is
nonincreasing and $t_2\ge t_{\mathrm{align}}$.

\smallskip\noindent(h) \emph{Below the prior.} At $t_2$, with $h'=h^\star-\frac{1}{2}\Delta$ and $\tanh z\ge\frac{z}{2}$ on $[0,1]$,
\begin{alignat*}{2}
 \Delta&=g\tanh(\tfrac{|a|}{2})\ge\frac{\initscale^2h^\star}{64}\cdot\frac{a_{\mathrm{sc}}}{32}=\Delta_2,\qquad \Delta\le\initscale &\qquad& \why{(e), (g), (a)}\\
 \tilde{\mathcal L}(h',\Delta)&=(1-p)\ell(h^\star)+p\,\ell(-h^\star+\Delta)\le H_p-p(1-p)\Delta+\tfrac{p}{8}\Delta^2 && \why{Taylor, $\ell''\le\tfrac{1}{4}$}\\
 &\le H_p-\tfrac{1}{2}p(1-p)\Delta && \why{$\Delta\le4(1-p)$}\\
 \tilde{\mathcal L}(h,\Delta)&\le\tilde{\mathcal L}(h',\Delta)+\Bigl|\pd{\mathcal L}{c}\Bigr|\,|h-h'| && \why{convexity in $h$}\\
 |h-h'|&\le\frac{1}{\kappa_0}\Bigl(\Bigl|\pd{\mathcal L}{c}\Bigr|+\tfrac{1}{8}|\Delta|\Bigr)+\tfrac{1}{2}|\Delta|\le\frac{38(1+h^\star)\initscale^6}{\kappa_0} && \why{(b), (c), (a)}\\
 \Bigl|\pd{\mathcal L}{c}\Bigr|\,|h-h'|&\le\frac{874(1+h^\star)\initscale^{12}}{\kappa_0}\le\tfrac{1}{4}p(1-p)\Delta_2 && \why{(c), \cref{eq:J-eps}}
\end{alignat*}
using $\ell(h^\star)=-\log(1-p)$, $\ell(-h^\star)=-\log p$ and $\ell'(-h^\star)=-(1-p)$. Hence
$\mathcal L(t_2)=\tilde{\mathcal L}(h,\Delta)\le H_p-\frac{1}{2}p(1-p)\Delta+\frac{1}{4}p(1-p)\Delta_2\le H_p-\frac{1}{4}p(1-p)\Delta_2$.
\end{proof}

\subsection{Phase III and the theorem}

\paragraph{Auxiliary constant correction.} For $|a|\ge a_{\mathrm{sc}}/8$, monotonicity gives\[|\tanh(a/2)|\ge\tanh(a_{\mathrm{sc}}/16)\ge a_{\mathrm{sc}}/32,\]because $a_{\mathrm{sc}}/16\le1$ and $\tanh x\ge x/2$ on $[0,1]$. Consequently the valid squared lower bound is $a_{\mathrm{sc}}^2/1024$. The manuscript's $1/256$ is slightly too large. The manuscript's conservative descent denominator $12288=3\cdot4\cdot1024$ already covers this correction, so its auxiliary time constants and the theorem's exponents can be retained.

\begin{lemma}[Descent]
\label{lem:J-descent}
Let $t_3=\inf\{t\ge t_2:\mathcal L(t)\le \mathcal L^\star\}$. Then
\[
 t_3-t_2\le\frac{6\cdot10^5d_{\mathrm h}}{p_{\min}^2\initscale^{12}},
\]
every prediction of $s_t$, $t\ge2$, has positive margin for all $t\ge t_3$, and $\mathcal L(t)\le\frac{10^5d_{\mathrm h}}{\initscale^{12}(t-t_3)}$ for $t>t_3$.
\end{lemma}

\begin{proof}
For $t\ge t_2$,
\begin{alignat*}{2}
 \tanh^2(\frac{a}{2})&\ge\tfrac{1}{1024}a_{\mathrm{sc}}^2 &\qquad& \why{$|a|\ge a_{\mathrm{sc}}/8$: \cref{lem:J-phase2,lem:J-align}}\\
 \ddt{\mathcal L}&\le-\tfrac{1}{3}\initscale^2v^2\tanh^2(\frac{a}{2})\,(\mathrm{err}_0^2+\mathrm{err}_1^2) && \why{\cref{lem:J-gradient}}\\
 &\le-\frac{\initscale^{12}}{12288\,d_{\mathrm h}}\,(\mathrm{err}_0^2+\mathrm{err}_1^2) && \why{$\initscale^2v^2\ge\tfrac{1}{4}\initscale^4$}
\end{alignat*}
While $\mathcal L>\mathcal L^\star$,
\begin{alignat*}{2}
 \ddt{\mathcal L}&\le-\frac{\initscale^{12}p_{\min}^2}{860160\,d_{\mathrm h}} &\qquad& \why{\cref{lem:J-levels}(d)}\\
 t_3-t_2&\le\frac{860160\log2\cdot d_{\mathrm h}}{p_{\min}^2\initscale^{12}} && \why{$\mathcal L(t_2)<\log2$}
\end{alignat*}
For $t\ge t_3$, $\mathcal L<p_{\min}\log2$, so all margins are positive (\cref{lem:J-levels}(b)), and
\begin{alignat*}{2}
 \ddt{\mathcal L}&\le-\frac{\initscale^{12}}{98304\,d_{\mathrm h}}\,\mathcal L^2 &\qquad& \why{\cref{lem:J-levels}(c)}\\
 \frac{d}{dt}\frac{1}{\mathcal L}&\ge\frac{\initscale^{12}}{98304\,d_{\mathrm h}} && 
\end{alignat*}
\end{proof}

\begin{proof}[Proof of \cref{thm:joint}]
\emph{Reduction.} By the reflection, $p<\frac{1}{2}$; then $p_{\min}=p$.

\emph{Margins and convergence.} \Cref{lem:J-prior,lem:J-phase2,lem:J-descent} give positive margins for all $t\ge t_3$ and $\Lstep\to0$, with
\[
 t_3\le\underbrace{\frac{2h^\star}{\kappa_0\initscale^3}}_{\text{Phase I}}+\underbrace{\frac{5\log\frac{1}{\initscale}}{\kappa_0\initscale^2}+T_{\mathrm{align}}}_{\text{Phase II}}
 +\underbrace{\frac{6\cdot10^5d_{\mathrm h}}{p_{\min}^2\initscale^{12}}}_{\text{Phase III}} .
\]

\emph{The stated bound.} The alignment term is the first term of the theorem:
\[
 T_{\mathrm{align}}=\frac{2048\,e\sqrt{d_{\mathrm h}}}{p(1-p)\,h^\star\mathrm{gap}\,\initscale^4}.
\]
The Phase I and relaxation terms are at most a constant times the Phase III term, because
\[
 h^\star\le\frac{1}{p_{\min}},\qquad p(1-p)\ge\frac{p_{\min}}{2},\qquad \frac{h^\star}{\kappa_0}\le\frac{2e}{p_{\min}^2},\qquad
 \initscale^{-3},\ \log\frac{1}{\initscale}\,\initscale^{-2}\le d_{\mathrm h}\initscale^{-12} .
\]
If $q(0)k(0)<0$, then $T_{\mathrm{align}}=0$ and $\gamma=1$, so $\mathrm{gap}$ does not enter.

\emph{Probability.} Let $X=\frac{|q(0)|}{\initscale}$ and $Y=\frac{|k(0)|}{\initscale}$ be independent and uniform on $[\frac{1}{2},2]$. Then $X^2$ has density
$\frac{1}{3\sqrt z}\le\frac{2}{3}$ on $[\frac{1}{4},4]$, so
\[
 \Pr\bigl(|D|<\mathrm{gap}\,\initscale^2\bigr)=\E\,\Pr\bigl(|X^2-Y^2|<\mathrm{gap}\bigm|Y\bigr)\le2\mathrm{gap}\cdot\tfrac{2}{3}=\tfrac{4}{3}\mathrm{gap} .
\]
\end{proof}

\paragraph{Why the gap condition.}\label{rem:J-balance} If $|q(0)|=|k(0)|$ and $q(0)k(0)>0$, then $D=0$, $S=\frac{2\sqrt{d_{\mathrm h}}|a|}{\initscale^2}$, and
\[
 \ddt{a}=-f(t)\,a,\qquad f=\frac{2\initscale^2}{\sqrt{d_{\mathrm h}}}\operatorname{sign}(a)\,g\,\sigma'(a)\,\mathrm{err} .
\]
Then $a$ keeps its sign, so alignment requires $g$ to change sign, and the crossing argument of \cref{lem:J-phase2}(d) fails. The condition $|D|\ge\mathrm{gap}\,\initscale^2$ excludes this case.



\subsection{The optimization comparison in the same architecture}
\begin{corollary}[Outcome training from the same initialization]
\label{cor:joint-outcome}
Take the architecture and initialization of \cref{thm:joint} for any $K\ge2$, with any fixed mask at the answer position, and $\mixrate$ as in
\cref{eq:lambda}. Gradient flow on $\Lout$, training $(\mathbf U,\mathbf W,q,k)$ or any larger set of blocks (the relative bias included) under any block
schedule, obeys the plateau bound of \cref{thm:outcome-end} for $\mixrate>0$ and $\initscale$ in the range of that theorem.
\end{corollary}

\begin{proof}
After division by $\initscale$, every map and embedding block has operator norm at most $2\sqrt K$, the count embeddings have singular values
$1$, every MLP bias has $|\tilde b_m|=1$, and the readout is zero: this is \cref{ass:init}. A fixed mask keeps every attention row a softmax over
its visible sites (\cref{ass:arch}) and only lowers the numbers of scores and sites. Hence \cref{prop:score-degree}, \cref{prop:network}(i)--(iii),
\cref{lem:score-sum,lem:end-smooth} and \cref{thm:outcome-end} apply unchanged for any schedule, and \cref{lem:bias-uniform} covers a trained bias.
\end{proof}


\paragraph{Resulting separation.} With $K=2$, the same mask, the same zero-readout seed, and the same trained variables $(U,W,q,k)$, take $\varepsilon\le\min(\varepsilon^\star,c\lambda,\varepsilon_0)$. Outcome flow obeys \cref{thm:outcome-end}, whereas step flow obeys \cref{thm:joint}. At fixed $p\ne1/2$ and fixed nonzero gap, the step bound is polynomial in $\varepsilon^{-1}$ and $d_{\mathrm h}$ and has no further dependence on $n$. The outcome bound grows as $\lambda^{-n}$. These bounds produce an actual ordering of the two guaranteed time scales whenever $c\lambda^{-n}$ exceeds the polynomial step upper bound. If $d_{\mathrm h}$ and $\varepsilon^{-1}$ grow polynomially with $n$, this ordering holds for sufficiently large $n$.

\paragraph{Scope of this rewrite.} The joint convergence result remains restricted to $K=2$, $p\ne1/2$, the fixed two-source gate mask, fixed directions and value factor, and the stated query--key imbalance condition for misaligned seeds. The step-credit theorem still gives an upper bound and a nonzero polynomial coefficient, rather than a quantitative high-probability lower bound at arbitrary random initialization. The outcome theorem bounds expected correct-answer probability and loss, rather than greedy-decoding accuracy. No conclusion has been strengthened by suppressing any of these restrictions.

\end{document}
